\documentclass[11pt]{article}
\usepackage[margin=0.9in]{geometry}
\usepackage[T1]{fontenc}
\usepackage{lmodern}
\usepackage{microtype}
\usepackage{amsmath,amssymb,amsthm,mathtools}
\usepackage{booktabs,longtable,tabularx,array,multirow}
\usepackage{enumitem}
\usepackage{xcolor}
\usepackage{listings}
\usepackage{fancyhdr}
\usepackage{hyperref}
\usepackage[nameinlink,noabbrev]{cleveref}
\usepackage{url}
\hypersetup{
  colorlinks=true,
  linkcolor=blue!45!black,
  citecolor=green!35!black,
  urlcolor=blue!55!black,
  pdftitle={The Meijer G-function in intuitionistic primitive recursive arithmetic},
  pdfauthor={Research report prepared for Vladimir Reshetnikov},
  pdfsubject={Constructive quantitative proofs and Lean formalization plan}
}
\pagestyle{fancy}
\fancyhf{}
\lhead{Meijer $G$ in iPRA}
\rhead{16 November 2026}
\cfoot{\thepage}
\setlength{\headheight}{14pt}
\setlength{\parindent}{0pt}
\setlength{\parskip}{0.52em}
\setlist{nosep,leftmargin=2em}
\allowdisplaybreaks
\emergencystretch=2em

\newtheorem{theorem}{Theorem}[section]
\newtheorem{lemma}[theorem]{Lemma}
\newtheorem{proposition}[theorem]{Proposition}
\newtheorem{corollary}[theorem]{Corollary}
\newtheorem{definition}[theorem]{Definition}
\newtheorem{remark}[theorem]{Remark}
\newtheorem{contract}[theorem]{Finite contract}
\newtheorem{audit}[theorem]{Source audit}
\newtheorem{roadmap}[theorem]{Formalization step}
\newtheorem{example}[theorem]{Example}

\newcommand{\N}{\mathbb N}
\newcommand{\Z}{\mathbb Z}
\newcommand{\Q}{\mathbb Q}
\newcommand{\R}{\mathbb R}
\newcommand{\C}{\mathbb C}
\newcommand{\dd}{\,\mathrm d}
\newcommand{\ii}{\mathrm i}
\newcommand{\e}{\mathrm e}
\newcommand{\Log}{\operatorname{Log}}
\newcommand{\Arg}{\operatorname{Arg}}
\newcommand{\Res}{\operatorname*{Res}}
\newcommand{\BallEq}{\operatorname{BallEq}}
\newcommand{\Safe}{\operatorname{Safe}}
\newcommand{\Adm}{\operatorname{Adm}}
\newcommand{\Eval}{\operatorname{Eval}}
\newcommand{\Hyp}{\operatorname{Hyp}}
\newcommand{\thetaop}{\vartheta}
\newcommand{\code}[1]{\texttt{#1}}
\newcommand{\statusP}{\textbf{P}}
\newcommand{\statusQ}{\textbf{Q}}
\newcommand{\statusC}{\textbf{C}}
\newcommand{\statusR}{\textbf{R}}
\newcommand{\statusX}{\textbf{X}}
\newcommand{\G}[5]{G_{#1,#2}^{#3,#4}\!\left(#5\right)}
\newcommand{\Garr}[6]{G_{#1,#2}^{#3,#4}\!\left(\left.\begin{matrix}#5\\#6\end{matrix}\right|z\right)}
\newcommand{\poch}[2]{\left(#1\right)_{#2}}
\newcommand{\abs}[1]{\left|#1\right|}
\newcommand{\norm}[1]{\left\lVert#1\right\rVert}

\lstset{
  basicstyle=\ttfamily\small,
  columns=fullflexible,
  breaklines=true,
  frame=single,
  backgroundcolor=\color{black!2},
  showstringspaces=false
}

\title{\textbf{The Meijer $G$-function in intuitionistic primitive recursive arithmetic}\\
\large Mellin--Barnes certificates, residue normal forms, special functions, and a Lean formalization roadmap}
\author{Research report prepared for Vladimir Reshetnikov}
\date{16 November 2026}

\begin{document}
\maketitle

\begin{abstract}
The Meijer $G$-function is normally introduced by an improper complex contour
integral and then used as a universal container for hypergeometric functions,
Bessel functions, incomplete gamma functions, transform kernels, and numerous
parameter identities.  None of the phrases ``arbitrary complex number'',
``choose a contour'', or ``take the analytic continuation'' is primitive in
intuitionistic primitive recursive arithmetic (iPRA).  We therefore replace
each such phrase by finite data: rational complex approximants, pole-clearance
certificates, a branch certificate for $\Log z$, primitive-recursive tail and
quadrature bounds, and a requested dyadic accuracy.

The central result is a constructive equivalence between three presentations:
a certified Mellin--Barnes integral, a residue expansion in generalized
hypergeometric series, and the solution of the Meijer differential equation
selected by explicit local data.  In the nonconfluent case the residue
expansion is the most economical iPRA definition.  Confluent poles are handled
by a finite Laurent-jet algebra and yield the expected powers of $\Log z$.
From this engine we prove, in corrected quantitative form, the definition and
three contour regimes, Slater's formulas, the differential equation,
nonresonant fundamental systems, hypergeometric conversion and polynomial
cases, parameter cancellation, power shift, inversion, multiplication,
derivative and contiguous relations, scaling series, Mellin/Euler/convolution
and Laplace formulas, and the elementary and higher-function special cases
listed by Wikipedia and MathWorld.  The broad Narain, generalized Laplace, and Meijer transform statements are
proved as finite-window theorems under explicit transform certificates; the
Wimp family is reduced to the same certified-kernel principle after expanding
its auxiliary kernel, but its long normalization and maximal test class are not
developed as a separate theorem.  The web pages omit the function-space and
convergence hypotheses needed for an unconditional inversion theorem.

The report carefully distinguishes a human constructive proof, executable
finite regression checks, a Lean host-level scaffold, and a completed
object-iPRA derivation.  The first three are supplied.  The final endpoint must
have type $[]\vdash_{\mathrm{PRAAx}}\varphi$ in ProveIt's checked
\code{ConciseIPRA} syntax and remains the stated formalization work; no analytic
kernel proof is claimed here.
\end{abstract}

\tableofcontents
\newpage

\section{Scope, source audit, and exact status}

\subsection{What the two source pages actually assert}

The Wikipedia page gives one definition, three contour types, convergence
conditions, a differential equation, two residue expansions, conversion to and
from generalized hypergeometric functions, a polynomial repair formula,
parameter identities, derivatives, recurrences, scaling expansions, integral
formulas, transform kernels, and a long table of special cases.  MathWorld gives
the two sign conventions for the Mellin--Barnes integral, the four-list Wolfram
Language convention, a generalized fourth parameter $r$, and five special
values.  We treat a displayed equation, its domain clause, and its convergence
clause as separate claims.  This is important: several displayed identities
are correct only after adding hypotheses that the prose suppresses.

The machine-readable ledger \code{claims/coverage.csv} contains the complete
itemization.  Its status codes are:
\begin{center}
\begin{tabularx}{0.94\textwidth}{@{}cX@{}}
\toprule
\statusP & proved by a finite algebraic, coefficient, residue, or tail argument developed in this report;\\
\statusQ & proved under an explicit quantitative admissibility, Fubini, or transform certificate;\\
\statusC & source statement corrected or restricted (branch, boundary, resonance, or convergence);\\
\statusR & reduced to a reusable certificate theorem, but a broad ``all suitable functions'' assertion is not made;\\
\statusX & literal source wording is undefined or cannot be true on the advertised whole domain.\\
\bottomrule
\end{tabularx}
\end{center}
A status describes the human mathematical argument, not a completed Lean
object proof.

\subsection{Main constructive theorem}

Fix integers $0\le m\le q$ and $0\le n\le p$, finite parameter lists
$\mathbf a=(a_1,\ldots,a_p)$ and $\mathbf b=(b_1,\ldots,b_q)$, a branch of
$\Log$, and an argument $z\ne0$.  Define
\begin{equation}\label{eq:phi-minus}
 \Phi(s)=
 \frac{\prod_{j=1}^{m}\Gamma(b_j-s)
       \prod_{j=1}^{n}\Gamma(1-a_j+s)}
      {\prod_{j=m+1}^{q}\Gamma(1-b_j+s)
       \prod_{j=n+1}^{p}\Gamma(a_j-s)}.
\end{equation}
The traditional notation is
\begin{equation}\label{eq:main-definition}
 G_{p,q}^{m,n}\!\left(\left.\begin{matrix}\mathbf a\\\mathbf b\end{matrix}\right|z\right)
 =\frac1{2\pi\ii}\int_L \Phi(s)z^s\dd s.
\end{equation}
In iPRA, \eqref{eq:main-definition} is not a primitive definition.  A datum
must also supply a rational polygonal or vertical contour, pole-clearance and
branch certificates, and primitive-recursive error bounds.

\begin{theorem}[Certified equivalence of presentations]\label{thm:master}
For every fixed coded parameter datum satisfying one of the admissibility
contracts in \cref{sec:definitions,sec:convergence}, and every requested
precision $k$, there is a primitive-recursive procedure producing a rational
complex ball of radius at most $2^{-k}$ which contains each of the following
quantities whenever that presentation is admissible:
\begin{enumerate}[label=(\roman*)]
\item the truncated Mellin--Barnes integral with certified tails;
\item the right- or left-residue normal form, including confluent Laurent jets;
\item the local solution of the Meijer differential equation selected by the
      same Frobenius data.
\end{enumerate}
Whenever two presentations are simultaneously admissible on the same branch
domain, their two balls overlap.  In the nonconfluent right-loop regime the
normal form is the finite sum of generalized hypergeometric series in
\cref{thm:slater-right}; in the left-loop regime it is
\cref{thm:slater-left}.  All algorithms and all error schedules are primitive
recursive.  The proof uses only finite sums, products, bounded searches,
quantifier-free induction, and previously specified finite Gamma and residue
contracts.
\end{theorem}

The proof occupies \cref{sec:hypergeom,sec:definitions,sec:convergence}.  Its
value is logical as well as computational: all later identities can be proved
on whichever presentation makes the claim finite.  Parameter identities are
proved at the integrand level; local analytic identities are proved
coefficientwise; integral transforms are proved by finite Fubini estimates.
No unrestricted identity theorem is needed.

\subsection{Three levels that must not be conflated}

The accompanying files establish three different things.
\begin{enumerate}
\item The present report gives a human-readable constructive proof with explicit
      finite contracts and error reductions.
\item The Python programs perform exact-rational recurrence checks and
      high-precision numerical regression tests.  They are evidence against
      transcription errors, not proofs of universal statements.
\item The Lean file gives executable finite definitions and target theorem
      signatures.  It intentionally contains clearly marked placeholders for
      the analytic contracts and does not masquerade as an object-iPRA proof.
\end{enumerate}
The final ProveIt endpoint is a term of
\[
 []\;\vdash[\mathrm{PRAAx}]\;\varphi.
\]
A theorem about Lean's host type \code{Complex} or \code{Rat} is not that
endpoint.  Once such a derivation is built, ProveIt's checked theorem
\code{pra\_iff\_native} transfers it to the concise native scalar calculus.

\subsection{Corrections to literal web-page wording}

\begin{audit}[Fundamental systems]
The displayed $p\le q$ and $p\ge q$ lists are fundamental systems only under a
nonresonance hypothesis, for example $b_i-b_j\notin\Z$ in the first list and
$a_i-a_j\notin\Z$ in the second, together with nonzero leading coefficients.
At resonance, logarithmic Frobenius solutions or parameter derivatives replace
some listed functions.  The differential equation remains valid.
\end{audit}

\begin{audit}[``For all $x$'']
The formulas for $\log(1+x)$, inverse trigonometric functions, powers, and
Lerch transcendent require a branch.  The Heaviside formulas are pointwise
statements on the positive real axis away from $x=1$ unless a convention at the
jump is supplied.  The Lerch series has a pole or divergence at many advertised
points.  We prove branch-qualified or finite-observation versions rather than
the literal universal quantifier.
\end{audit}

\begin{audit}[Transform inversions]
The pages use ``suitable function'' and omit substantial convergence
conditions.  In weak arithmetic this cannot be left implicit.  We prove a
finite transform theorem from named endpoint, decay, quadrature, and Fubini
certificates.  Narain, Wimp, and the generalized Laplace formulas then become
instances.  We do not claim an unspecified maximal function-space theorem.
\end{audit}

\section{What complex analysis means in iPRA}\label{sec:ipra}

\subsection{The object theory}

The repository snapshot used here is commit
\code{85afffc93b2abbc1972fdfc98244052cbfbda1b3}.  Its ordinary presentation of
iPRA has first-order intuitionistic logic with equality, a symbol for every
primitive-recursive function, the defining equations for those symbols,
successor axioms, and induction for quantifier-free formulas.  In Lean these
are the types \code{PR}, \code{PTm}, and \code{PRAAx}; the induction constructor
is \code{PRA.qfInd}.  Consequently, a proof may use an externally fixed
primitive-recursive approximation algorithm, but it may not quantify over all
algorithms by invoking a universal evaluator.

This distinction controls the formulation of every theorem below.  A statement
such as ``for every complex parameter $a$'' is read as a schema: fix a
primitive-recursive name for $a$ and the certificates demanded by the theorem;
then iPRA proves the resulting first-order sentence.  For rational parameters
the package also gives a single uniform coded algorithm.

\subsection{Rational and complex codes}

A rational is represented by a sign and two natural numbers $(u,v)$ with
$v>0$; normalization is optional.  Equality is the cross-multiplied equation
$u_1v_2=u_2v_1$, so no gcd computation is logically required.  A rational
complex number is a pair of rational codes.  We use the $\ell^1$ magnitude
\[
 \norm{x+\ii y}_1=\abs{x}+\abs{y},
\]
because all its upper-bound relations reduce to integer inequalities.  It is
submultiplicative up to the exact estimate
$\norm{zw}_1\le\norm z_1\norm w_1$.

\begin{definition}[Dyadic ball equality]
For rational complex codes $u,v$ and $k\in\N$,
\[
 \BallEq_k(u,v)\quad:\Longleftrightarrow\quad
 \norm{u-v}_1\le 2^{-k}.
\]
After clearing positive denominators, this is quantifier-free arithmetic.
\end{definition}

A complex name is an externally fixed primitive-recursive function
$u:\N\to\Q(\ii)$ satisfying
\begin{equation}\label{eq:name-coherence}
 \BallEq_k(u(k+2),u(k+1))
 \quad\text{and}\quad
 \BallEq_k(u(k+1),u(k)).
\end{equation}
After a harmless reindexing, the geometric series proves that all sufficiently
accurate approximants represent one observational value.  No completed Cauchy
real sort is added to the language.

\subsection{Apartness, branches, and safe division}

Division by $w$ is permitted only with an apartness witness
$\norm w_1\ge2^{-r}$.  A branch certificate for $z$ consists of:
\begin{enumerate}
\item an apartness witness $\norm z_1\ge2^{-r}$;
\item a rational angular interval $I$ of width below $2\pi$ containing the
      argument of every allowed approximation to $z$;
\item a rational clearance from the chosen branch cut;
\item primitive-recursive approximants to $\Log z$ with a modulus derived from
      those clearances.
\end{enumerate}
Then $z^a$ means $\exp(a\Log z)$ on that named branch.  All identities involving
powers use the same branch on both sides.  This removes a frequent hidden
classical error: $z^{a+b}=z^az^b$ is not a branch-free identity.

For the positive real special cases, a branch certificate is just a positive
rational lower bound, and the construction simplifies substantially.

\subsection{Finite analytic contracts}

\begin{contract}[Gamma]\label{con:gamma}
For a fixed compact rational box $K$ with rational distance $\eta>0$ from the
nonpositive integers, \code{GammaApprox} returns, from $(w,k)$, a rational
complex ball of radius $2^{-k}$ containing $\Gamma(w)$ uniformly for $w\in K$.
It also returns derivative jets of any externally fixed order and vertical-line
bounds
\begin{align}
 \abs{\Gamma(\sigma+\ii t)}
 &\le C(1+\abs t)^{\sigma-1/2}\e^{-\pi\abs t/2},\label{eq:gamma-upper}\\
 \abs{1/\Gamma(\sigma+\ii t)}
 &\le C'(1+\abs t)^{1/2-\sigma}\e^{\pi\abs t/2},\label{eq:recip-gamma-upper}
\end{align}
for $\abs t\ge T_0$, with rational $C,C',T_0$ computed from $K$ and $\eta$.
\end{contract}
The existing \code{gamma\_ipra} packet develops the finite Euler products and
Stirling bounds needed for this contract.  That packet itself explicitly does
not claim a completed analytic object derivation; our formalization plan treats
\cref{con:gamma} as a module interface to be discharged.

\begin{contract}[Power series]
A series $\sum c_r$ is certified by a primitive-recursive index $N(k)$ and an
upper bound $E(N)$ such that finite arithmetic proves
$\sum_{r\ge N}\norm{c_r}_1\le E(N)\le2^{-k}$.  Differentiated series carry the
same data for each externally fixed derivative order.
\end{contract}

\begin{contract}[Contour integral]
A contour is a rational polygonal path, or a vertical segment with rational
real part, together with: a lower distance from every pole and branch cut; a
bound on the integrand and a fixed number of derivatives on each segment; a
finite quadrature rule; and endpoint/tail errors.  The value at precision $k$
is a finite rational sum with total error at most $2^{-k}$.
\end{contract}

\begin{contract}[Improper real integral and Fubini]
An integral on $(0,\infty)$ is split at rational $0<\epsilon<1<R$.  The middle
piece has a finite quadrature certificate; the two endpoint pieces have
explicit majorants.  A change of integration order additionally supplies a
majorant for the discarded rectangles and an error budget proving that both
orders approximate the same finite double sum.
\end{contract}

All four contracts are arithmetical data.  For fixed analytic names, their
verification consists of bounded loops and quantifier-free induction.

\subsection{Formula shape in the object language}

A typical analytic theorem is not written inside iPRA as an equality of
abstract complex numbers.  Its precision-$k$ instance has the form
\begin{equation}\label{eq:object-shape}
 \forall x\le B(k)\;
   \bigl(\Adm(d,x,k)\to
   \BallEq_k(\Eval_{\rm left}(d,x,k+c),
             \Eval_{\rm right}(d,x,k+c))\bigr),
\end{equation}
where $c$ is a fixed guard-bit budget.  Every symbol in
\eqref{eq:object-shape} is primitive recursive.  Unbounded parameters are
handled by a schema in the external choices $d$ and $B$, not by a universal
program interpreter.

The finite triangle inequality gives transitivity of ball equality.  Thus a
proof graph may pass from a contour evaluator to a residue evaluator, then to a
hypergeometric evaluator, and finally to an elementary evaluator.  The total
precision budget is allocated backwards through that graph.  This is the
mechanism by which the many special cases become ordinary iPRA derivations.

\section{The finite hypergeometric engine}\label{sec:hypergeom}

The residue formulas reduce most of the theory to generalized hypergeometric
series.  We therefore prove their finite theory before mentioning contours.

\subsection{Coefficients and primitive recursion}

For complex parameters $\mathbf A=(A_1,\ldots,A_r)$ and
$\mathbf B=(B_1,\ldots,B_s)$ define
\begin{equation}\label{eq:hypergeom}
 {}_rF_s\!\left(\left.\begin{matrix}\mathbf A\\\mathbf B\end{matrix}\right|w\right)
 =\sum_{k=0}^{\infty}t_k,
 \qquad
 t_0=1,
 \qquad
 t_{k+1}=t_k\,w\,
 \frac{\prod_{i=1}^{r}(A_i+k)}{(k+1)\prod_{j=1}^{s}(B_j+k)}.
\end{equation}
For rational complex inputs, the partial sum and the pair
$(t_k,\sum_{j<k}t_j)$ are obtained by simultaneous primitive recursion.
Denominators are protected by a certificate that no $B_j$ is a nonpositive
integer, or by termination before the first zero when the series is a
polynomial.

\begin{lemma}[Rising-factorial form]\label{lem:rising}
For every $k$ before a denominator zero,
\[
 t_k=\frac{\prod_{i=1}^{r}\poch{A_i}{k}}
           {\prod_{j=1}^{s}\poch{B_j}{k}}\frac{w^k}{k!}.
\]
\end{lemma}
\begin{proof}
The assertion at $k=0$ is a finite ring identity.  Multiplication by the ratio
in \eqref{eq:hypergeom} gives the successor case, using
$\poch A{k+1}=\poch Ak(A+k)$ and $(k+1)!=k!(k+1)$.  This is precisely a
quantifier-free induction invariant after clearing denominators.
\end{proof}

\subsection{A uniform ratio-tail certificate}

Convergence is estimated with the ordinary complex modulus.  A rational
certificate $\abs w\le W$ is checked without a square root by
$(\Re w)^2+(\Im w)^2\le W^2$; center uncertainty is absorbed into $W$ by the
triangle inequality.  The $\ell^1$ norm remains the representation used for
output balls, and $\norm z_1\le2\abs z$ costs only one guard bit.

Let $M_i\ge\abs{A_i}$, and let $L_j(k)>0$ be rational lower bounds for
$\abs{B_j+k}$.  When $t_k\ne0$, the coefficient recurrence yields
\begin{equation}\label{eq:ratio-majorant}
 \frac{\abs{t_{k+1}}}{\abs{t_k}}
 \le W\,
 \frac{\prod_i(k+M_i)}{(k+1)\prod_jL_j(k)}.
\end{equation}
If $t_k=0$, every later term is zero under denominator apartness, so the same
tail conclusion is immediate.  If a rational $0<\rho<1$ and an index $N$ make
the right side at most $\rho$ for every $k\ge N$, then
\begin{equation}\label{eq:geometric-tail}
 \abs{\sum_{k=N}^{\infty}t_k}
 \le\frac{\abs{t_N}}{1-\rho}.
\end{equation}
The proof is finite induction followed by the finite geometric-sum identity;
for a requested precision one truncates the geometric sum after a further
explicit number of terms and bounds its last term.  Converting the final bound
to the $\ell^1$ ball uses the factor $2$ above.  No completed infinite sum is
needed internally.

\begin{theorem}[Primitive-recursive convergence]\label{thm:hyp-conv}
For fixed parameter names with denominator-apartness certificates:
\begin{enumerate}[label=(\alph*)]
\item if $r\le s$, then for every bounded $w$ a primitive-recursive search
      produces $N$ and $\rho=1/2$ satisfying \eqref{eq:ratio-majorant};
\item if $r=s+1$ and a squared-modulus certificate proves $\abs w\le R<1$,
      choose $R<\rho<1$; a primitive-recursive search produces $N$ satisfying
      the ratio bound;
\item if one upper parameter is $-h$ with $h\in\N$, then $t_k=0$ for $k>h$,
      so the series is a finite polynomial without any convergence principle.
\end{enumerate}
The resulting truncation modulus is primitive recursive in the requested
precision and all rational certificate data.
\end{theorem}
\begin{proof}
After multiplying by all positive denominator bounds, the desired ratio
inequality is a polynomial inequality in $k$.  In case (a) the denominator has
strictly larger degree; in case (b) the leading-coefficient ratio is bounded by
$R<\rho$.  A coarse explicit threshold is obtained by replacing every
$k+M_i$ by $(1+M)k$ and each eventually positive $L_j(k)$ by $k/2$.
The finitely many preceding indices are checked directly.  A bounded search up
to the resulting threshold is primitive recursive.  Equation
\eqref{eq:geometric-tail} supplies the modulus.  Case (c) follows from the
factor $-h+h=0$.
\end{proof}

\subsection{Termwise algebra and derivatives}

Uniform ratio bounds on a rational compact set persist after multiplying the
$k$th term by any fixed polynomial in $k$.  Therefore one may differentiate an
externally fixed number of times, multiply series, or shift parameters while
retaining a primitive-recursive tail bound.

\begin{lemma}[Coefficientwise identity principle]\label{lem:coeff-id}
Suppose two certified power or power-log series on the same branch have equal
coefficients by a quantifier-free arithmetic proof and have uniform tail
moduli.  Then their evaluations are ball-equal to every requested precision.
\end{lemma}
\begin{proof}
Given $k$, choose a common truncation index whose two tails are below
$2^{-k-2}$.  The finite partial sums are exactly equal after clearing
 denominators.  Two applications of the triangle inequality give error below
$2^{-k-1}$, which is stronger than required after guard-bit reindexing.
\end{proof}
This finite lemma replaces almost every appeal to the classical identity
theorem in the local part of the report.

\subsection{Terminating and singular parameter cases}

A lower parameter in $\{0,-1,-2,\ldots\}$ makes the naive recurrence undefined.
There are three constructive resolutions, each recorded in the datum:
\begin{enumerate}
\item an upper $-h$ terminates earlier, so all used denominators are nonzero;
\item a common factor is cancelled symbolically before evaluation;
\item a confluent limit is represented by finite parameter jets and a
      quantitative divided-difference bound.
\end{enumerate}
The polynomial formula on the source page is needed because the naive Gamma
prefactor in the one-$G$ representation may vanish exactly when the
hypergeometric series terminates.  We prove the repaired two-$G$ formula in
\cref{sec:hypergeom-conversion} by finite coefficient extraction.

\section{Meijer data and the three definitions}\label{sec:definitions}

\subsection{Four parameter blocks}

It is logically and computationally preferable to store the parameters in the
four blocks
\[
 A^-=(a_1,\ldots,a_n),\quad A^+=(a_{n+1},\ldots,a_p),\quad
 B^-=(b_1,\ldots,b_m),\quad B^+=(b_{m+1},\ldots,b_q).
\]
This is also the Wolfram Language convention
\code{MeijerG[\{\{A-\},\{A+\}\},\{\{B-\},\{B+\}\},z]}.
The separation condition for the standard contour is
\begin{equation}\label{eq:pole-separation}
 a_k-b_j\notin\{1,2,3,\ldots\}
 \qquad(1\le k\le n,\ 1\le j\le m),
\end{equation}
because the two pole families are
$s=b_j+\ell$ and $s=a_k-1-\ell$.
In iPRA, \eqref{eq:pole-separation} is replaced by a rational lower bound on
their distance over the finite contour window under consideration, plus an
arithmetic proof that the two arithmetic progressions never meet.

Define
\begin{equation}\label{eq:delta-nu}
 \delta=m+n-\frac{p+q}{2},
 \qquad
 \nu=\sum_{j=1}^{q}b_j-\sum_{j=1}^{p}a_j.
\end{equation}
These quantities govern vertical-line decay.

\subsection{Mellin--Barnes definition and sign convention}

The Wikipedia/MathWorld convention is \eqref{eq:main-definition}.  Setting
$u=-s$ gives the equally standard inverse-Mellin form
\begin{equation}\label{eq:plus-convention}
 \frac1{2\pi\ii}\int_{-L}
 \frac{\prod_{j=1}^{m}\Gamma(b_j+u)
       \prod_{j=1}^{n}\Gamma(1-a_j-u)}
      {\prod_{j=m+1}^{q}\Gamma(1-b_j-u)
       \prod_{j=n+1}^{p}\Gamma(a_j+u)}z^{-u}\dd u.
\end{equation}
The orientation reversal from $\dd s=-\dd u$ is cancelled by reversing the
endpoints.  This is a finite change-of-variable proof on every truncated
contour; the tail certificates are identical because $\abs u=\abs s$.
Thus MathWorld equations (1) and (2) are constructively equivalent.

\begin{definition}[Vertical datum]
A vertical datum consists of a rational $\sigma$ such that the line
$\Re s=\sigma$ separates the two pole families, a branch interval for
$\Log z$, and constants certifying the bound in
\cref{thm:vertical-convergence}.  Its precision-$k$ value is a rational
quadrature on $\sigma-\ii T(k)\le s\le\sigma+\ii T(k)$.
\end{definition}

\begin{definition}[Right-loop and left-loop data]
A right-loop datum encloses the poles $b_j+\ell$ and excludes the poles
$a_k-1-\ell$; a left-loop datum does the reverse.  At precision $k$ only
finitely many poles and a finite closing rectangle are used.  The datum gives
a remainder bound showing that the omitted side tends below $2^{-k}$.
\end{definition}

\subsection{Residue normal forms}

For nonconfluent right poles, define $G$ directly by the finite sum of
hypergeometric names in \cref{thm:slater-right}.  For nonconfluent left poles,
use \cref{thm:slater-left}.  These definitions are preferable in iPRA: they
reduce evaluation to primitive recursion and a ratio-tail certificate, and the
contour integral becomes a theorem rather than a foundational infinitary
object.

When poles coalesce, the correct normal form is a power-log series.  We now give
a finite algorithm that covers every multiplicity.

\begin{definition}[Laurent-jet residue algorithm]\label{def:jet}
Let $s_0$ be a pole of multiplicity $r$.  Put $u=s-s_0$.  For each Gamma or
reciprocal-Gamma factor, request its Laurent or Taylor jet through degree
$r-1$ with a rational remainder ball.  Expand
\[
 z^{s_0+u}=z^{s_0}\sum_{j=0}^{r-1}\frac{(u\Log z)^j}{j!}+O(u^r).
\]
Multiply the finite Laurent polynomials by bounded convolution and select the
coefficient of $u^{-1}$.  The result is
$z^{s_0}$ times a polynomial in $\Log z$ of degree at most $r-1$.
\end{definition}
Every step in \cref{def:jet} is primitive recursive for externally fixed $r$.
The remainder is bounded by finite Cauchy products of the jet remainders.  In a
pole sequence indexed by $\ell$, the additional polynomial factor in $\ell$
and $\Log z$ does not destroy the geometric or superfactorial tail proved in
\cref{thm:hyp-conv}.  Thus confluent cases such as integer-order $K_\nu$ and the
repeated-parameter Lerch formulas require no limiting choice principle.

\subsection{Generalized fourth parameter}

MathWorld's generalized form is
\begin{equation}\label{eq:r-general}
 G_{p,q}^{m,n}\!\left(\left.\begin{matrix}\mathbf a\\\mathbf b\end{matrix}
 \right|x;r\right)
 =\frac1{2\pi\ii}\int_L
 \Psi(s)x^{-s/r}\dd s,
\end{equation}
where $\Psi$ is the plus-sign kernel in \eqref{eq:plus-convention}.  For a fixed
nonzero rational $r$ and a branch of $\Log x$, this is covered by the same
vertical and residue machinery: replace $\Log x$ by $\Log x/r$ in every
exponential estimate.  The substitution $s=ru$ yields the finite identity
\begin{equation}\label{eq:r-substitution}
 G(x;r)=\frac{r}{2\pi\ii}\int_{L/r}\Psi(ru)x^{-u}\dd u.
\end{equation}
For integer or reciprocal-integer $r$, Gauss multiplication for Gamma converts
\eqref{eq:r-substitution} to an ordinary $G$ with expanded parameter lists and
an explicit power and constant.  The special $r=1/2$ cosine formula is proved
directly in \cref{thm:mathworld-specials}.

\subsection{Admissibility predicate}

A coded datum is \emph{admissible at precision $k$} when all of the following
finite checks pass:
\begin{enumerate}
\item list lengths and block indices are consistent;
\item every requested denominator and Gamma factor has an apartness witness;
\item the branch box avoids zero and the branch cut;
\item the contour window has rational pole clearance;
\item the supplied tail, jet, and quadrature bounds sum to at most $2^{-k}$;
\item every use of an improper real integral carries endpoint majorants;
\item every exchange of sums or integrals carries a finite Fubini budget.
\end{enumerate}
This predicate, denoted $\Adm(d,k)$, is primitive recursive.  All global-looking
analytic assertions in the sequel are abbreviations for the family of finite
implications $\Adm(d,k)\to\BallEq_k(\cdots)$.

\section{Convergence, contour comparison, and residues}\label{sec:convergence}

\subsection{The vertical-line estimate}

Let $s=\sigma+\ii t$ and let $\phi=\Arg z$ on the certified branch.  Applying
the two Gamma bounds \eqref{eq:gamma-upper}--\eqref{eq:recip-gamma-upper} to
each factor of \eqref{eq:phi-minus} gives a computable constant $C>0$ such
that
\begin{equation}\label{eq:master-vertical-bound}
 \abs{\Phi(\sigma+\ii t)z^{\sigma+\ii t}}
 \le C\abs z^\sigma(1+\abs t)^A
      \exp\!\bigl(-[\pi\delta-\abs\phi]\abs t\bigr),
\end{equation}
where
\begin{equation}\label{eq:A-exponent}
 A=\Re\nu+(p-q)\left(\sigma+\frac12\right).
\end{equation}
Indeed, each numerator Gamma contributes $\e^{-\pi\abs t/2}$ and each
reciprocal denominator Gamma contributes $\e^{\pi\abs t/2}$, leaving
$\e^{-\pi\delta\abs t}$.  The powers of $1+\abs t$ from all $a$-parameters and
$b$-parameters sum exactly to \eqref{eq:A-exponent}.  Finally,
$\abs{z^s}=\abs z^\sigma\e^{-t\phi}$, whose symmetric upper bound is
$\abs z^\sigma\e^{\abs\phi\abs t}$.

\begin{theorem}[Constructive vertical convergence]\label{thm:vertical-convergence}
Suppose the line $\Re s=\sigma$ has certified pole clearance.
\begin{enumerate}[label=(\alph*)]
\item If $\epsilon:=\pi\delta-\abs{\Arg z}>0$ with a rational lower bound,
      then the vertical integral converges with a primitive-recursive tail
      modulus.
\item On the boundary $\epsilon=0$, it converges whenever a rational
      $\eta>0$ certifies $A\le-1-\eta$; equivalently,
      $(q-p)(\sigma+1/2)\ge\Re\nu+1+\eta$.
\end{enumerate}
These are the quantitative versions of the first contour conditions on the
source page.
\end{theorem}
\begin{proof}
In case (a), put $M=\max(0,\lceil A\rceil)$.  For
$T\ge(4M/\epsilon)^2+1$, the elementary estimate
$\log(1+t)\le\sqrt t$ gives
$(1+t)^M\le\e^{\epsilon t/2}$.  Therefore the two tails are at most
\[
 \frac{4C\abs z^\sigma}{\epsilon}\e^{-\epsilon T/2}.
\]
A bounded search using the finite Taylor enclosure for $\exp$ finds a $T(k)$
for which this is below $2^{-k}$.  In case (b),
\[
 \int_T^\infty C(1+t)^{-1-\eta}\dd t
 \le \frac C\eta T^{-\eta}.
\]
Writing $\eta=r/s$ rational, choose $T$ by the integer inequality
$T^r\ge(2^kC/\eta)^s$.  Integer powering and bounded root search are primitive
recursive.  The finite central segment is covered by rational quadrature.
\end{proof}

\subsection{Right residues and Slater's formula}

Assume no two $b_1,\ldots,b_m$ differ by an integer.  The right poles are
$s=b_h+k$, $h=1,\ldots,m$, $k\ge0$.  The residue of
$\Gamma(b_h-s)$ with respect to $s$ is
$(-1)^{k+1}/k!$.  The right loop is negatively oriented, so its extra minus
sign changes this to $(-1)^k/k!$.  Apply
\begin{equation}\label{eq:gamma-shift-identities}
 \Gamma(x+k)=\Gamma(x)\poch{x}{k},
 \qquad
 \Gamma(x-k)=\frac{(-1)^k\Gamma(x)}{\poch{1-x}{k}}
\end{equation}
to every remaining factor.

\begin{theorem}[Right-residue expansion / Slater]\label{thm:slater-right}
If $p<q$, or if $p=q$ and $\abs z<1$ with rational margin, and the right poles
are nonconfluent, then
\begin{align}
&G_{p,q}^{m,n}\!\left(\left.\begin{matrix}\mathbf a\\\mathbf b\end{matrix}\right|z\right)
 =\sum_{h=1}^{m}C_h z^{b_h}
 {}_pF_{q-1}\!\left(
 \left.\begin{matrix}1+b_h-\mathbf a\\
 (1+b_h-\mathbf b)^{\!*}
 \end{matrix}\right|(-1)^{p-m-n}z\right),\label{eq:slater-right}\\
&C_h=
 \frac{\prod_{j=1,\,j\ne h}^{m}\Gamma(b_j-b_h)
       \prod_{j=1}^{n}\Gamma(1+b_h-a_j)}
      {\prod_{j=m+1}^{q}\Gamma(1+b_h-b_j)
       \prod_{j=n+1}^{p}\Gamma(a_j-b_h)}.
\end{align}
The star means that the lower parameter with index $h$ is omitted.  The equality
has a primitive-recursive modulus on every certified compact branch domain.
\end{theorem}
\begin{proof}
The residue calculation above gives the $k$th term.  Substituting
\eqref{eq:gamma-shift-identities} and collecting the signs yields precisely the
coefficient recurrence \eqref{eq:hypergeom}.  When $p<q$, the series on the
right has at least as many lower parameters as upper parameters and converges
for all bounded $z$ by \cref{thm:hyp-conv}; when $p=q$, it has radius one and
the certified margin supplies a geometric tail.  A finite right-closing
rectangle contains finitely many poles.  The residue theorem for that polygon
is a finite triangulation identity plus local Laurent-jet calculations.  Gamma
ratio bounds make the outer side smaller than the same hypergeometric tail.
Given $k$, choose the rectangle and series cutoffs so that both errors are below
$2^{-k-2}$; the finite residue sums agree exactly.  Ball transitivity proves
\eqref{eq:slater-right}.
\end{proof}

\begin{corollary}[Empty right family]\label{cor:mzero}
For the right-loop definition,
$G_{p,q}^{0,n}=0$ when $p<q$, or when $p=q$ and $\abs z<1$.
\end{corollary}
\begin{proof}
Every finite right contour contains no poles, so its integral is zero by the
finite Cauchy theorem.  The closing-side error tends below every requested
dyadic tolerance.  Equivalently, the sum in \eqref{eq:slater-right} is empty.
\end{proof}

\subsection{Left residues}

The poles enclosed by the left loop are $s=a_h-1-k$ for $h\le n$.  The same
calculation, or inversion from \cref{thm:inversion}, gives the reciprocal
formula.

\begin{theorem}[Left-residue expansion]\label{thm:slater-left}
If $p>q$, or if $p=q$ and $\abs z>1$ with rational margin, and the left poles
are nonconfluent, then
\begin{align}
&G_{p,q}^{m,n}\!\left(\left.\begin{matrix}\mathbf a\\\mathbf b\end{matrix}\right|z\right)
 =\sum_{h=1}^{n}D_h z^{a_h-1}
 {}_qF_{p-1}\!\left(
 \left.\begin{matrix}1-a_h+\mathbf b\\
 (1-a_h+\mathbf a)^{\!*}
 \end{matrix}\right|(-1)^{q-m-n}z^{-1}\right),\label{eq:slater-left}\\
&D_h=
 \frac{\prod_{j=1,\,j\ne h}^{n}\Gamma(a_h-a_j)
       \prod_{j=1}^{m}\Gamma(1-a_h+b_j)}
      {\prod_{j=n+1}^{p}\Gamma(1-a_h+a_j)
       \prod_{j=m+1}^{q}\Gamma(a_h-b_j)}.
\end{align}
\end{theorem}
The proof is the reflected version of \cref{thm:slater-right}.  The argument of
the hypergeometric series has norm below one when $p=q$ because it is
proportional to $z^{-1}$.

\begin{corollary}[Empty left family]\label{cor:nzero}
For the left-loop definition, $G_{p,q}^{m,0}=0$ when $p>q$, or when $p=q$ and
$\abs z>1$.
\end{corollary}

\begin{theorem}[Balanced unit-circle boundary]\label{thm:balanced-unit-circle}
Suppose $p=q$, $\abs z=1$, and a rational margin $\eta>0$ certifies
$\Re\nu\le-1-\eta$.  Then both the right-loop and left-loop residue
presentations converge whenever their pole-separation data are admissible.
For every fixed confluent multiplicity they have a primitive-recursive
$p$-series tail modulus as well.
\end{theorem}
\begin{proof}
Consider one right Slater summand.  Its $k$th coefficient, apart from its fixed
prefactor and the unit-modulus phase, is a quotient of Pochhammer symbols.
Writing each Pochhammer symbol as a Gamma quotient and using the fixed-strip
Stirling bounds gives computable $C,N$ such that
\begin{equation}\label{eq:unit-circle-coeff}
 \abs{t_{h,k}}\le C(k+1)^{\Re\nu+\eta/2}
 \le C(k+1)^{-1-\eta/2}\qquad(k\ge N).
\end{equation}
The exponent follows by direct cancellation:
\[
 \sum_{i=1}^{p}\Re(1+b_h-a_i)
 -\sum_{j\ne h}\Re(1+b_h-b_j)-1=\Re\nu.
\]
The Stirling statement can alternatively be checked directly from the squared
term-ratio polynomial: its leading coefficient is $1$ and its next coefficient
is $2\Re\nu$, so bounded search finds the same $N$ with the prescribed margin.
For $a=\eta/2$,
\[
 \sum_{k\ge N}(k+1)^{-1-a}\le a^{-1}N^{-a};
\]
when $a=r/s$ is rational, the required cutoff is obtained from an integer
inequality after raising to the power $s$.  This is a primitive-recursive
$p$-series certificate.

A pole group of fixed multiplicity contributes an additional fixed polynomial
in $\Log z$, harmonic sums, and $\log(k+1)$.  Since bounded search proves
$(1+\log(k+1))^d\le(k+1)^{\eta/4}$ eventually for every fixed $d$, the same
argument leaves exponent at most $-1-\eta/4$.  The left expansion is identical
after inversion, because $\abs{z^{-1}}=1$ and the same $\nu$ occurs.  The
closing-side estimates are bounded by these residue tails, completing the two
loop integrals.
\end{proof}

\subsection{Confluent poles}

If $b_j-b_h\in\Z$, several right-pole sequences overlap.  The simple constants
$C_h$ then contain divergent Gamma factors, but the original contour integral
need not diverge.  Group the overlapping poles and apply the jet algorithm
\cref{def:jet}.  A multiplicity-$r$ group contributes
\begin{equation}\label{eq:power-log-normal}
 z^{\beta+k}\sum_{j=0}^{r-1}c_{k,j}(\Log z)^j.
\end{equation}
The coefficients $c_{k,j}$ are finite expressions in Gamma jets, factorials,
and harmonic sums.  Their ratios are bounded by a fixed polynomial in $k$
times the nonconfluent majorant, so the same convergence domains hold.  The
left case is identical.  Thus the ``no confluent poles'' hypothesis belongs to
the short closed form \eqref{eq:slater-right}, not to the existence of the
$G$-function.

\subsection{Agreement of contour types and analyticity}

\begin{theorem}[Contour agreement]\label{thm:contour-agreement}
Whenever two of the vertical, right-loop, and left-loop presentations satisfy
their quantitative hypotheses on the same branch domain, they are ball-equal
to every precision.
\end{theorem}
\begin{proof}
Truncate both contours and join them by a rational rectangle avoiding every
pole.  Finite triangulation cancels interior edges.  The remaining difference
is the sum of residues in the strip plus the outer sides.  If no poles lie
between the contours, the residue sum is empty.  Otherwise both presentations
are compared with the same finite residue normal form.  The vertical estimate
\eqref{eq:master-vertical-bound} and the right/left series tails make all outer
pieces arbitrarily small with primitive-recursive cutoffs.
\end{proof}

\begin{theorem}[Quantitative analyticity]\label{thm:analyticity}
On a compact rational domain that avoids $z=0$, the chosen branch cut, and in
the balanced case the singular point selected by the differential equation,
the certified $G$-name has Taylor coefficients and a primitive-recursive
uniform remainder modulus.  Every externally fixed derivative order is
therefore represented by a primitive-recursive name.
\end{theorem}
\begin{proof}
In a residue domain, differentiate the uniformly convergent power-log series;
polynomial factors in the summation index preserve the ratio tail.  In a
vertical domain, differentiate $z^s$ an externally fixed number of times and
use \eqref{eq:master-vertical-bound} with a larger polynomial exponent.  A
finite Cauchy estimate on a rational circle converts derivative bounds into a
Taylor remainder.  All precision choices are bounded searches.
\end{proof}

\section{Differential equation and local solution bases}\label{sec:ode}

Put $\thetaop=z\frac{\dd}{\dd z}$.  The differential equation listed on
Wikipedia is
\begin{equation}\label{eq:meijer-ode}
 \left[(-1)^{p-m-n}z\prod_{j=1}^{p}(\thetaop-a_j+1)
       -\prod_{j=1}^{q}(\thetaop-b_j)\right]G(z)=0.
\end{equation}

\subsection{The kernel shift identity}

The proof is a one-line Gamma recurrence once the signs are tracked correctly.
From $\Gamma(w+1)=w\Gamma(w)$,
\begin{equation}\label{eq:phi-shift}
 \frac{\Phi(s)}{\Phi(s-1)}
 =(-1)^{p-m-n}
   \frac{\prod_{j=1}^{p}(s-a_j)}
        {\prod_{j=1}^{q}(s-b_j)}.
\end{equation}
The exponent can also be written $p+m-n$; the two differ by the even number
$2m$.

\begin{theorem}[Meijer differential equation]\label{thm:ode}
Every certified presentation of $G_{p,q}^{m,n}$ satisfies
\eqref{eq:meijer-ode} in the quantitative derivative sense of
\cref{thm:analyticity}.
\end{theorem}
\begin{proof}
On a finite vertical contour, $\thetaop$ multiplies $z^s$ by $s$.  Hence the
second product in \eqref{eq:meijer-ode} has integrand
$\prod_j(s-b_j)\Phi(s)z^s$.  In the first product, shift $s$ to $s-1$ after the
extra factor $z$; its integrand becomes
$(-1)^{p-m-n}\prod_j(s-a_j)\Phi(s-1)z^s$, equal by
\eqref{eq:phi-shift}.  The two finite contours differ only by horizontal
connector segments and residues crossed by the shift.  Choosing a common
admissible strip, or performing the argument on the residue normal form, makes
these corrections cancel.  Tail and derivative errors are allocated below
$2^{-k-2}$.  In the nonconfluent series presentation, the same result follows
more directly: inserting $z^{b_h+k}$ turns \eqref{eq:meijer-ode} into the
coefficient recurrence obtained from \eqref{eq:hypergeom}.  For confluent
power-log terms, $\thetaop$ acts on the finite polynomial in $\Log z$, and the
Laurent-jet identity gives the same cancellation.
\end{proof}

\subsection{Order and singularities}

Expanding the products in $\thetaop$ shows that the equation has order
$\max(p,q)$ except at a degenerate leading coefficient.  At $z=0$ it is a
regular singular equation when $p\le q$ in the standard local direction; at
$z=\infty$ the inverted equation is analogous.  When $p=q$, the leading
coefficient vanishes at
\begin{equation}\label{eq:balanced-singularity}
 z=(-1)^{p-m-n},
\end{equation}
which explains the additional singular point mentioned on the source page.
These are finite polynomial calculations in the coefficients of
\eqref{eq:meijer-ode}.

\subsection{Fundamental systems and the missing hypothesis}

For $p\le q$, place $b_h$ first in the bottom list and consider
\begin{equation}\label{eq:right-basis}
 Y_h(z)=G_{p,q}^{1,p}\!\left(
 \left.\begin{matrix}a_1,\ldots,a_p\\
 b_h,b_1,\ldots,\widehat{b_h},\ldots,b_q
 \end{matrix}\right|(-1)^{p-m-n+1}z\right),
 \quad h=1,\ldots,q.
\end{equation}
The right-residue formula gives
\begin{equation}\label{eq:right-leading}
 Y_h(z)=c_hz^{b_h}(1+O(z)),\qquad c_h\ne0,
\end{equation}
under Gamma apartness.  Apply $\dd^{r}/\dd z^{r}$ for
$r=0,\ldots,q-1$.  The leading Wronskian is
\begin{equation}\label{eq:wronskian}
 \left(\prod_{h=1}^{q}c_h\right)
 z^{\sum_hb_h-q(q-1)/2}
 \prod_{1\le i<j\le q}(b_j-b_i),
\end{equation}
up to a factor $1+O(z)$.

\begin{theorem}[Corrected fundamental-system theorem]\label{thm:fundamental}
Suppose $p\le q$, the $b_h$ are pairwise distinct modulo integers, and every
leading Gamma factor in \eqref{eq:right-leading} has a quantitative apartness
witness.  Then the $q$ functions \eqref{eq:right-basis} form a fundamental
system on a sufficiently small certified punctured sector.  Dually, for
$p\ge q$, the $p$ functions obtained by placing each $a_h$ first and using the
left-residue presentation form a fundamental system when the $a_h$ are
pairwise distinct modulo integers.
\end{theorem}
\begin{proof}
The Vandermonde factor in \eqref{eq:wronskian} has a rational lower bound by
hypothesis.  Uniform tail bounds make the $O(z)$ correction smaller than half
the leading term on an explicitly searched rational radius.  Therefore the
Wronskian is apart from zero there.  Each function solves the same equation by
\cref{thm:ode}; an order-$q$ linear equation cannot have more than $q$
independent local solutions.  The usual linear-algebra proof is finite: solve
the initial-value system by Cramer's rule using the Wronskian apartness witness.
The left case follows after $z\mapsto z^{-1}$.
\end{proof}

\begin{remark}[Resonance]
If two exponents differ by an integer, the Vandermonde argument fails and the
listed $G$-functions may be dependent.  The jet construction supplies the
missing logarithmic Frobenius solutions.  Thus the source statement is correct
generically but not literally without a resonance clause.
\end{remark}

\subsection{Uniqueness without an unrestricted identity theorem}

Many later special cases are established by showing that two named functions
solve the same differential equation and have the same local coefficients.  In
this report that phrase abbreviates a finite argument: for every requested
precision choose a rational disk, compute enough coefficients, bound both
remainders, and use a finite recurrence to show coefficient equality.  If one
uses initial values instead, the Wronskian lower bound gives a finite
Gr\"onwall/Picard estimate.  No quantification over arbitrary open sets is
required.

\section{Hypergeometric conversion and polynomial cases}\label{sec:hypergeom-conversion}

\subsection{Every convergent generalized hypergeometric series is a \texorpdfstring{$G$}{G}}

Put one bottom parameter $0$ in the right block and transform the others by
$a_i\mapsto1-a_i$, $b_j\mapsto1-b_j$.  In
\cref{thm:slater-right}, there is only the pole family $b_1=0$.  Its prefactor
is
\[
 \frac{\prod_i\Gamma(a_i)}{\prod_j\Gamma(b_j)},
\]
and its hypergeometric argument is $z$.  Therefore:

\begin{theorem}[Hypergeometric-to-Meijer conversion]\label{thm:hyp-to-g}
Assume no upper parameter is a nonpositive integer, no denominator parameter
is a nonpositive integer, and the series is in its convergence domain.  Common
upper/lower factors may first be cancelled.  Then
\begin{align}
{}_pF_q\!\left(\left.\begin{matrix}\mathbf a\\\mathbf b\end{matrix}\right|z\right)
&=\frac{\prod_{j=1}^{q}\Gamma(b_j)}
        {\prod_{i=1}^{p}\Gamma(a_i)}
G_{p,q+1}^{1,p}\!\left(
 \left.\begin{matrix}1-\mathbf a\\0,1-\mathbf b\end{matrix}\right|-z\right),
\label{eq:hyp-g-first}\\
&=\frac{\prod_{j=1}^{q}\Gamma(b_j)}
        {\prod_{i=1}^{p}\Gamma(a_i)}
G_{q+1,p}^{p,1}\!\left(
 \left.\begin{matrix}1,\mathbf b\\\mathbf a\end{matrix}\right|-z^{-1}\right).
\label{eq:hyp-g-second}
\end{align}
The first equality holds for every finite $z$ when $p\le q$, and for
$\abs z<1$ when $p=q+1$.  The second is the first followed by inversion and is
read on the correspondingly transported branch.
\end{theorem}
\begin{proof}
The calculation preceding the theorem turns \eqref{eq:slater-right} into
$G=(\prod\Gamma(a_i)/\prod\Gamma(b_j)){}_pF_q$.  All coefficients are equal
exactly by \cref{lem:rising}; the tails have the common modulus from
\cref{thm:hyp-conv}.  Equation \eqref{eq:hyp-g-second} follows from
\cref{thm:inversion}.  Gamma prefactors are evaluated only with pole-apartness
witnesses.
\end{proof}

\subsection{Analytic continuation and the case \texorpdfstring{$p>q+1$}{p > q + 1}}

When $p=q+1$, the series has radius one.  Formula
\eqref{eq:hyp-g-first} supplies a continuation along every certified path that
avoids the selected cut, because overlapping disks are compared by finite
coefficient equality and explicit reexpansion bounds.  The familiar principal
continuation has a cut from $1$ to $+\infty$ after the argument $-z$ and branch
conventions are translated.

For $p>q+1$, the generic hypergeometric power series diverges for every
$z\ne0$.  A $G$-contour nevertheless defines sectorial solutions of the same
formal differential equation.  This is a useful extension, but not a unique
branch-free value: different admissible contours can encode different Stokes
solutions.  The precise constructive statement is:

\begin{proposition}[Sectorial extension]
Fix a contour type, a branch sector, and all vertical or residue tail
certificates.  Then \eqref{eq:hyp-g-first}, read as a definition rather than an
equality to a convergent series, gives a primitive-recursive name solving the
formal generalized hypergeometric equation and having the prescribed
asymptotic/Frobenius data.  Two such names agree where their contours are
certifiably deformable into one another.
\end{proposition}
This qualifies the page's phrase ``the natural extension'' without weakening
the computational content.

\subsection{Terminating series}

If an upper parameter equals $-h$, the Gamma prefactor in
\eqref{eq:hyp-g-first} contains $1/\Gamma(-h)=0$, while the contour parameters
simultaneously become resonant.  The limiting product can be finite and
nonzero; multiplying a formal infinity by zero is not a proof.  The source page
repairs the problem with a sum of two $G$-functions.

Let $0\le m\le q$, $0\le n\le p$, and write
\[
 R_{m,n}(\mathbf a,\mathbf b)=
 \frac{\prod_{j=n+1}^{p}\Gamma(1-a_j)
       \prod_{j=m+1}^{q}\Gamma(b_j)}
      {\prod_{j=1}^{n}\Gamma(a_j)
       \prod_{j=1}^{m}\Gamma(1-b_j)}.
\]

\begin{theorem}[Polynomial repair formula]\label{thm:polynomial-repair}
For $h\in\N$, wherever the displayed prefactors are interpreted after the
permitted cancellations,
\begin{align}
&{}_{p+1}F_q\!\left(\left.\begin{matrix}-h,\mathbf a\\\mathbf b\end{matrix}\right|z\right)
=h!R_{m,n}(\mathbf a,\mathbf b)\,[\mathcal G_+(z)+(-1)^h\mathcal G_-(z)],
\label{eq:poly-repair}\\
&\mathcal G_+(z)=G_{p+1,q+1}^{m+1,n}\!\left(
 \left.\begin{matrix}1-\mathbf a,h+1\\0,1-\mathbf b\end{matrix}\right|
 (-1)^{p-m-n}z\right),\\
&\mathcal G_-(z)=G_{p+1,q+1}^{m,n+1}\!\left(
 \left.\begin{matrix}h+1,1-\mathbf a\\1-\mathbf b,0\end{matrix}\right|
 (-1)^{p-m-n}z\right).
\end{align}
This is a finite iPRA identity: the left side has degree $h$, and the right side
has the same coefficients through degree $h$ and zero coefficients afterward.
\end{theorem}
\begin{proof}
Expand $\mathcal G_+$ by the right residues and $\mathcal G_-$ by the left
residues, using Laurent jets if a chosen $m,n$ creates a collision.  Apply
\eqref{eq:gamma-shift-identities} and the reflection formula
$\Gamma(u)\Gamma(1-u)=\pi/\sin\pi u$.  For the coefficient of $z^r$, all
parameter-dependent constants reduce to
\[
 \frac{\poch{-h}{r}\prod_i\poch{a_i}{r}}
      {r!\prod_j\poch{b_j}{r}}.
\]
For $0\le r\le h$ this is the polynomial coefficient.  For $r>h$ the factor
$\poch{-h}{r}$ is zero; in the residue presentation this zero is the exact
cancellation of the two complementary terms.  The calculation is a finite
ring identity after multiplying by the relevant Gamma products.  Since every
coefficient after $h$ is zero, no infinite tail remains on the combined right
side.  Quantifier-free induction on $r\le h$ proves the identity in iPRA.
\end{proof}

\begin{remark}
The freedom to choose $m$ and $n$ is useful, not cosmetic.  One chooses blocks
so that no individual prefactor asks for an undefined Gamma value; if a
collision is unavoidable, the jet evaluator computes the finite combined
limit before division.
\end{remark}

\section{Parameter algebra: cancellation, shifts, inversion, multiplication}\label{sec:algebra}

All identities in this section are first proved on a finite truncated contour.
They are therefore algebraic identities of rational approximants plus explicit
tail transport.  Alternatively, they can be checked coefficientwise in the
residue normal forms.

\subsection{Cancellation}

If $a_1=b_q$ with $a_1$ in the numerator $a$-block and $b_q$ in the denominator
$b$-block, the factors
$\Gamma(1-a_1+s)$ and $\Gamma(1-b_q+s)$ cancel.  If $a_p=b_1$ with $a_p$ in the
denominator $a$-block and $b_1$ in the numerator $b$-block, the factors
$\Gamma(a_p-s)$ and $\Gamma(b_1-s)$ cancel.

\begin{theorem}[Order reduction]\label{thm:cancellation}
With parameters placed in the indicated blocks,
\begin{align}
G_{p,q}^{m,n}\!\left(
 \left.\begin{matrix}a_1,a_2,\ldots,a_p\\b_1,\ldots,b_{q-1},a_1\end{matrix}\right|z\right)
&=G_{p-1,q-1}^{m,n-1}\!\left(
 \left.\begin{matrix}a_2,\ldots,a_p\\b_1,\ldots,b_{q-1}\end{matrix}\right|z\right),\label{eq:cancel-one}\\
G_{p,q}^{m,n}\!\left(
 \left.\begin{matrix}a_1,\ldots,a_{p-1},b_1\\b_1,b_2,\ldots,b_q\end{matrix}\right|z\right)
&=G_{p-1,q-1}^{m-1,n}\!\left(
 \left.\begin{matrix}a_1,\ldots,a_{p-1}\\b_2,\ldots,b_q\end{matrix}\right|z\right).
\label{eq:cancel-two}
\end{align}
\end{theorem}
\begin{proof}
The two Mellin kernels are literally equal after cancelling a common factor.
The same contour and every finite error certificate are reusable.  If a
numerical evaluator avoids symbolic cancellation, cross-multiply by the common
Gamma enclosure and use its apartness witness.
\end{proof}

\subsection{Power shift}

\begin{theorem}[Power shift]\label{thm:power-shift}
On one fixed branch of $\Log z$,
\begin{equation}\label{eq:power-shift}
 z^\rho G_{p,q}^{m,n}\!\left(
 \left.\begin{matrix}\mathbf a\\\mathbf b\end{matrix}\right|z\right)
 =G_{p,q}^{m,n}\!\left(
 \left.\begin{matrix}\mathbf a+\rho\\\mathbf b+\rho\end{matrix}\right|z\right).
\end{equation}
\end{theorem}
\begin{proof}
In the right side put $u=s-\rho$.  Every Gamma argument becomes the
corresponding unshifted argument, while $z^s=z^\rho z^u$.  Translate the finite
contour by $\rho$; pole-clearance certificates translate with it.  On residue
series, every exponent $b_h+k$ increases by $\rho$ and all parameter
differences are unchanged.  The branch certificate proves the power law.
\end{proof}

\subsection{Integer interchange identities}

Let $d=\alpha'-\alpha\in\Z$.  The reflection formula gives
\begin{equation}\label{eq:reflection-ratio}
 \frac{\Gamma(1-\alpha+s)/\Gamma(\alpha'-s)}
      {\Gamma(1-\alpha'+s)/\Gamma(\alpha-s)}
 =\frac{\sin\pi(\alpha'-s)}{\sin\pi(\alpha-s)}=(-1)^d.
\end{equation}
This proves the top interchange identity.  The same argument with $b$-factors
proves the bottom identity.

\begin{theorem}[Integer interchanges]\label{thm:interchanges}
Whenever the two sides have common admissible contours,
\begin{align}
&G_{p+2,q}^{m,n+1}\!\left(
 \left.\begin{matrix}\alpha,\mathbf a,\alpha'\\\mathbf b\end{matrix}\right|z\right)
 =(-1)^{\alpha'-\alpha}
 G_{p+2,q}^{m,n+1}\!\left(
 \left.\begin{matrix}\alpha',\mathbf a,\alpha\\\mathbf b\end{matrix}\right|z\right),\label{eq:top-swap}\\
&G_{p,q+2}^{m+1,n}\!\left(
 \left.\begin{matrix}\mathbf a\\\beta,\mathbf b,\beta'\end{matrix}\right|z\right)
 =(-1)^{\beta'-\beta}
 G_{p,q+2}^{m+1,n}\!\left(
 \left.\begin{matrix}\mathbf a\\\beta',\mathbf b,\beta\end{matrix}\right|z\right).
\label{eq:bottom-swap}
\end{align}
If $d=\beta-\alpha\in\N$, then
\begin{equation}\label{eq:block-transfer}
G_{p+1,q+1}^{m,n+1}\!\left(
 \left.\begin{matrix}\alpha,\mathbf a\\\mathbf b,\beta\end{matrix}\right|z\right)
=(-1)^d
G_{p+1,q+1}^{m+1,n}\!\left(
 \left.\begin{matrix}\mathbf a,\alpha\\\beta,\mathbf b\end{matrix}\right|z\right).
\end{equation}
\end{theorem}
\begin{proof}
Equations \eqref{eq:top-swap} and \eqref{eq:bottom-swap} follow from
\eqref{eq:reflection-ratio}.  For \eqref{eq:block-transfer}, repeated use of
$\Gamma(w+1)=w\Gamma(w)$ shows that the quotient of the two affected kernel
factors is the constant $(-1)^d$.  All sine and Gamma divisions are performed
only after the common factors have been cross-multiplied; the final identity
extends through removable parameter collisions by the jet evaluator.
\end{proof}

\subsection{Inversion}

\begin{theorem}[Inversion]\label{thm:inversion}
With compatible branches,
\begin{equation}\label{eq:inversion}
G_{p,q}^{m,n}\!\left(
 \left.\begin{matrix}\mathbf a\\\mathbf b\end{matrix}\right|z\right)
=G_{q,p}^{n,m}\!\left(
 \left.\begin{matrix}1-\mathbf b\\1-\mathbf a\end{matrix}\right|z^{-1}\right).
\end{equation}
\end{theorem}
\begin{proof}
Substitute $u=-s$ in the truncated contour integral.  Then
$z^s=(z^{-1})^u$, the factors $\Gamma(b_j-s)$ become
$\Gamma(1-(1-b_j)+u)$, and the factors $\Gamma(1-a_j+s)$ become
$\Gamma((1-a_j)-u)$ in the complementary block.  Orientation signs cancel as
in \eqref{eq:plus-convention}.  A right loop becomes a left loop, and the
branch certificate transports $\Log(z^{-1})=-\Log z$.  The tail bounds are
unchanged.
\end{proof}

\subsection{Gauss multiplication}

For $h\ge1$, write
\[
 \Delta_h(c)=\left(\frac c h,\frac{c+1}h,\ldots,\frac{c+h-1}h\right).
\]
Apply the constructive Gamma multiplication theorem
\begin{equation}\label{eq:gauss-mult}
 \Gamma(hw)=(2\pi)^{(1-h)/2}h^{hw-1/2}
 \prod_{r=0}^{h-1}\Gamma\!\left(w+\frac rh\right)
\end{equation}
to every factor after the substitution $s=hu$.

\begin{theorem}[Meijer multiplication theorem]\label{thm:multiplication}
For $h\in\N_{>0}$ and $z\ne0$ on a fixed branch,
\begin{align}
&G_{p,q}^{m,n}\!\left(
 \left.\begin{matrix}\mathbf a\\\mathbf b\end{matrix}\right|z\right)\notag\\
&\quad=
\frac{h^{1+\nu+(p-q)/2}}{(2\pi)^{(h-1)\delta}}
G_{hp,hq}^{hm,hn}\!\left(
 \left.\begin{matrix}\Delta_h(a_1),\ldots,\Delta_h(a_p)\\
                       \Delta_h(b_1),\ldots,\Delta_h(b_q)
 \end{matrix}\right|
 \frac{z^h}{h^{h(q-p)}}\right).
\label{eq:meijer-multiplication}
\end{align}
\end{theorem}
\begin{proof}
Put $s=hu$ in the left integral.  Expand every Gamma factor by
\eqref{eq:gauss-mult}.  The $(2\pi)$ powers combine to
$(2\pi)^{-(h-1)\delta}$; the constant powers of $h$ combine with $\dd s=h\dd u$
to $h^{1+\nu+(p-q)/2}$.  The $u$-dependent powers of $h$ join $z^{hu}$ to give
$(z^h/h^{h(q-p)})^u$.  The remaining Gamma factors are exactly the expanded
parameter lists.  This is an identity on each finite contour; the Gamma
multiplication packet transports the error bounds.  All exponents of $h$ are
rational powers on the positive real branch and therefore have direct finite
names.
\end{proof}

\section{Derivatives, recurrences, and scaling series}\label{sec:derivatives}

\subsection{Four first-order derivative formulas}

The identity $\frac{\dd}{\dd z}z^s=sz^{s-1}$ combines with a single Gamma
recurrence.  For example,
\[
 (s-a_1+1)\Gamma(1-a_1+s)=\Gamma(2-a_1+s),
\]
which lowers $a_1$ by one.  Depending on the block, the sign is reversed.

\begin{theorem}[First-order formulas]\label{thm:first-derivatives}
On every certified differentiability domain,
\begin{align}
\frac{\dd}{\dd z}\!
 \left[z^{1-a_1}G(\mathbf a;\mathbf b\mid z)\right]
&=z^{-a_1}G(a_1-1,a_2,\ldots;\mathbf b\mid z),&&n\ge1,
\label{eq:d-a1}\\
\frac{\dd}{\dd z}\!
 \left[z^{1-a_p}G(\mathbf a;\mathbf b\mid z)\right]
&=-z^{-a_p}G(a_1,\ldots,a_p-1;\mathbf b\mid z),&&n<p,
\label{eq:d-ap}\\
\frac{\dd}{\dd z}\!
 \left[z^{-b_1}G(\mathbf a;\mathbf b\mid z)\right]
&=-z^{-1-b_1}G(\mathbf a;b_1+1,b_2,\ldots\mid z),&&m\ge1,
\label{eq:d-b1}\\
\frac{\dd}{\dd z}\!
 \left[z^{-b_q}G(\mathbf a;\mathbf b\mid z)\right]
&=z^{-1-b_q}G(\mathbf a;b_1,\ldots,b_q+1\mid z),&&m<q.
\label{eq:d-bq}
\end{align}
\end{theorem}
\begin{proof}
Differentiate a finite contour integral or finite residue partial sum.  In each
case, collect the factor $s+1-a_i$ or $s-b_j$ with the affected Gamma or
reciprocal-Gamma factor and use $\Gamma(w+1)=w\Gamma(w)$.  Uniform derivative
tails follow from \cref{thm:analyticity}.  The sign in
\eqref{eq:d-ap} and \eqref{eq:d-b1} is the sign obtained from
$a_p-s=-(s-a_p)$ or $b_1-s=-(s-b_1)$.
\end{proof}

Expanding \eqref{eq:d-a1} gives the useful Euler-derivative identity
\begin{equation}\label{eq:theta-shift}
 \thetaop G(\mathbf a;\mathbf b\mid z)
 =G(a_1-1,a_2,\ldots;\mathbf b\mid z)+(a_1-1)G(\mathbf a;\mathbf b\mid z).
\end{equation}
The other three analogous identities are obtained in the same finite way.

\subsection{Arbitrary derivative order}

For $h\in\N$, the falling factorial
$s^{\underline h}=s(s-1)\cdots(s-h+1)$ is represented by adding a top parameter
$0$ and a bottom parameter $h$, because
\begin{equation}\label{eq:falling-gamma}
 \frac{\Gamma(1+s)}{\Gamma(1-h+s)}=s^{\underline h}.
\end{equation}

\begin{theorem}[Higher derivatives]\label{thm:higher-derivatives}
For each fixed $h$,
\begin{align}
z^h\frac{\dd^h}{\dd z^h}G_{p,q}^{m,n}(\mathbf a;\mathbf b\mid z)
&=G_{p+1,q+1}^{m,n+1}(0,\mathbf a;\mathbf b,h\mid z)\label{eq:high-d-one}\\
&=(-1)^hG_{p+1,q+1}^{m+1,n}(\mathbf a,0;h,\mathbf b\mid z).
\label{eq:high-d-two}
\end{align}
For the inverse argument,
\begin{align}
z^h\frac{\dd^h}{\dd z^h}G_{p,q}^{m,n}(\mathbf a;\mathbf b\mid z^{-1})
&=G_{p+1,q+1}^{m+1,n}(\mathbf a,1-h;1,\mathbf b\mid z^{-1})\notag\\
&=(-1)^hG_{p+1,q+1}^{m,n+1}(1-h,\mathbf a;\mathbf b,1\mid z^{-1}).
\label{eq:high-d-inverse}
\end{align}
\end{theorem}
\begin{proof}
The $h$th derivative of $z^s$ is $s^{\underline h}z^{s-h}$.
Equation \eqref{eq:falling-gamma} inserts the two new parameters and proves
\eqref{eq:high-d-one}.  The block-transfer identity
\eqref{eq:block-transfer} with difference $h$ proves
\eqref{eq:high-d-two}.  Apply inversion and the finite chain rule to obtain
\eqref{eq:high-d-inverse}.  Since $h$ is externally fixed, all derivative
error bounds remain primitive recursive.
\end{proof}

\subsection{Contiguous recurrences}

Equating two versions of $\thetaop G$ gives the source page's recurrences.
They are therefore not new analytic theorems; they are finite linear
combinations of \cref{thm:first-derivatives}.

\begin{theorem}[Four displayed contiguous relations]\label{thm:contiguous}
Using the shorthand $G=G(\mathbf a;\mathbf b\mid z)$,
\begin{align}
(a_p-a_1)G
&=G(a_1-1,a_2,\ldots;\mathbf b)
 +G(a_1,\ldots,a_p-1;\mathbf b),&&1\le n<p,
\label{eq:contig-aa}\\
(b_1-b_q)G
&=G(\mathbf a;b_1+1,b_2,\ldots)
 +G(\mathbf a;b_1,\ldots,b_q+1),&&1\le m<q,
\label{eq:contig-bb}\\
(b_1-a_1+1)G
&=G(a_1-1,a_2,\ldots;\mathbf b)
 +G(\mathbf a;b_1+1,b_2,\ldots),&&m,n\ge1,
\label{eq:contig-ab}\\
(a_p-b_q-1)G
&=G(a_1,\ldots,a_p-1;\mathbf b)
 +G(\mathbf a;b_1,\ldots,b_q+1),&&n<p,\ m<q.
\label{eq:contig-ba}
\end{align}
\end{theorem}
\begin{proof}
Expand the left derivatives in \eqref{eq:d-a1}--\eqref{eq:d-bq}, solve each for
$\thetaop G$, and subtract the appropriate pair.  This uses only distributivity
and integer/rational coefficient arithmetic, so it compiles to a short object
proof once the four derivative lemmas are available.
\end{proof}

The two omitted diagonal recurrences mentioned by the page follow by adding or
subtracting \eqref{eq:contig-aa}--\eqref{eq:contig-ba}; the checker can generate
them automatically.

\subsection{Scaling/Taylor formulas}

The four source expansions are Taylor series in the scale parameter $w$ around
$w=1$.  We prove one; the others follow from the corresponding first-order
formula.  Put
\[
 H(w)=w^{-b_1}G(\mathbf a;\mathbf b\mid wz).
\]
Using \eqref{eq:d-b1} and the chain rule gives
\[
 H^{(h)}(w)=(-1)^h w^{-b_1-h}
 G(\mathbf a;b_1+h,b_2,\ldots\mid wz),
 \qquad
 \frac{H^{(h)}(1)}{h!}
 =\frac{(-1)^h}{h!}G(\mathbf a;b_1+h,b_2,\ldots\mid z).
\]
Rewriting $(-1)^h(w-1)^h=(1-w)^h$ gives the first displayed
series.  The other three are identical after replacing the first-order shift
identity.

\begin{theorem}[Scaling series]\label{thm:scaling-series}
On every rational disk in the $w$-plane for which the common branch and a
uniform derivative bound are certified,
\begin{align}
G(\mathbf a;\mathbf b\mid wz)
&=w^{b_1}\sum_{h\ge0}\frac{(1-w)^h}{h!}
  G(\mathbf a;b_1+h,b_2,\ldots\mid z),&&m\ge1,
\label{eq:scale-b1}\\
G(\mathbf a;\mathbf b\mid wz)
&=w^{b_q}\sum_{h\ge0}\frac{(w-1)^h}{h!}
  G(\mathbf a;b_1,\ldots,b_q+h\mid z),&&m<q,
\label{eq:scale-bq}\\
G(\mathbf a;\mathbf b\mid z/w)
&=w^{1-a_1}\sum_{h\ge0}\frac{(1-w)^h}{h!}
  G(a_1-h,a_2,\ldots;\mathbf b\mid z),&&n\ge1,
\label{eq:scale-a1}\\
G(\mathbf a;\mathbf b\mid z/w)
&=w^{1-a_p}\sum_{h\ge0}\frac{(w-1)^h}{h!}
  G(a_1,\ldots,a_p-h;\mathbf b\mid z),&&n<p.
\label{eq:scale-ap}
\end{align}
\end{theorem}
\begin{proof}
Taylor's formula with an integral or Cauchy remainder is applied on a certified
rational disk.  The derivatives are replaced by shifted $G$-functions using
\cref{thm:first-derivatives,thm:higher-derivatives}.  A uniform Cauchy bound
$M/R^h$ gives the tail
$M(\abs{w-1}/R)^{N+1}/(1-\abs{w-1}/R)$, with a primitive-recursive cutoff.
This also explains why the page correctly warns that the radius depends on
$z$ and on the function being expanded.
\end{proof}

\section{Mellin, Euler, convolution, Laplace, and transform pairs}\label{sec:transforms}

The source page explicitly omits many convergence restrictions.  We retain the
formulas but replace ``whenever the integral exists'' by checkable endpoint,
strip, branch, and Fubini certificates.

\subsection{A finite Mellin-pair theorem}

Write the plus-sign kernel as
\begin{equation}\label{eq:M-kernel}
 M(s)=
 \frac{\prod_{j=1}^{m}\Gamma(b_j+s)
       \prod_{j=1}^{n}\Gamma(1-a_j-s)}
      {\prod_{j=m+1}^{q}\Gamma(1-b_j-s)
       \prod_{j=n+1}^{p}\Gamma(a_j+s)}.
\end{equation}
The expected pair is
\begin{equation}\label{eq:mellin-pair}
 G(\eta x)=\frac1{2\pi\ii}\int_{c-\ii\infty}^{c+\ii\infty}
 M(s)(\eta x)^{-s}\dd s,
 \qquad
 \int_0^\infty x^{u-1}G(\eta x)\dd x=\eta^{-u}M(u).
\end{equation}
A raw exchange of these two infinite integrals would insert a Dirac delta and
is not a valid elementary Fubini proof.  The constructive proof uses finite
Fourier inversion after $x=\e^y$.

\begin{definition}[Mellin certificate]\label{def:mellin-cert}
A Mellin-pair certificate at $c$ gives:
\begin{enumerate}
\item endpoint majorants for $f(y)=\e^{cy}G(\eta\e^y)$ and its required finite
      differences as $y\to\pm\infty$;
\item vertical decay and quadrature bounds for $M(c+\ii t)$;
\item a rational modulus of continuity for $f$ on every requested finite
      interval;
\item an error schedule for the Fej\'er kernels
\[
 K_T(v)=\frac1{2\pi T}\left(\frac{\sin(Tv/2)}{\sin(v/2)}\right)^2
\]
(or the continuous triangular Fourier analogue).
\end{enumerate}
\end{definition}
All occurrences of trigonometric functions in the finite kernel have certified
power-series approximants; its positivity and normalization reduce to finite
identities and rational error bounds.

\begin{theorem}[Certified Mellin transform]\label{thm:mellin}
Assume a Mellin certificate and a rational strip
\begin{equation}\label{eq:mellin-strip}
 -\min_{1\le j\le m}\Re b_j+\epsilon
 \le\Re u\le
 1-\max_{1\le j\le n}\Re a_j-\epsilon
\end{equation}
with $\epsilon>0$ (with the absent side omitted when $m=0$ or $n=0$).  Then
\begin{equation}\label{eq:mellin-result}
 \int_0^\infty x^{u-1}G_{p,q}^{m,n}(\eta x)\dd x
 =\eta^{-u}
 \frac{\prod_{j=1}^{m}\Gamma(b_j+u)
       \prod_{j=1}^{n}\Gamma(1-a_j-u)}
      {\prod_{j=m+1}^{q}\Gamma(1-b_j-u)
       \prod_{j=n+1}^{p}\Gamma(a_j+u)}.
\end{equation}
The equality is a family of finite ball equalities with a primitive-recursive
modulus.
\end{theorem}
\begin{proof}
After $x=\e^y$, the Mellin transform on the line $\Re u=c$ is the Fourier
transform of $f(y)$.  Truncate $y$ and $t$ to rational intervals, use finite
quadrature, and replace the sharp cutoff by the triangular multiplier.  The
resulting double sum may be reordered exactly.  The finite Fourier kernel
approximates the identity by the supplied modulus of continuity, while endpoint
and frequency tails use the certificate.  Moving $u$ within the strip is
justified by multiplying $f(y)$ by $\e^{(\Re u-c)y}$; the margin $\epsilon$
preserves both tails.  The Gamma ratio is $M(u)$ by definition.  Every error is
assigned a dyadic budget before the finite computation.
\end{proof}

\subsection{Euler--Beta integrals}

Insert the finite contour representation of $G(zx)$ and use
\begin{equation}\label{eq:beta-inner}
 \int_0^1 x^{s-\alpha}(1-x)^{\alpha-\beta-1}\dd x
 =\frac{\Gamma(1-\alpha+s)\Gamma(\alpha-\beta)}
        {\Gamma(1-\beta+s)}.
\end{equation}
The new Gamma quotient is exactly one new top numerator parameter $\alpha$ and
one new bottom denominator parameter $\beta$.

\begin{theorem}[Two Euler-type integrals]\label{thm:euler-integrals}
Assume $\Re(\alpha-\beta)>0$, compatible branches, and an endpoint/Fubini
certificate.  Then
\begin{align}
\int_0^1x^{-\alpha}(1-x)^{\alpha-\beta-1}
G_{p,q}^{m,n}(zx)\dd x
&=\Gamma(\alpha-\beta)
G_{p+1,q+1}^{m,n+1}\!\left(
 \left.\begin{matrix}\alpha,\mathbf a\\\mathbf b,\beta\end{matrix}\right|z\right),
\label{eq:euler-zero-one}\\
\int_1^\infty x^{-\alpha}(x-1)^{\alpha-\beta-1}
G_{p,q}^{m,n}(zx)\dd x
&=\Gamma(\alpha-\beta)
G_{p+1,q+1}^{m+1,n}\!\left(
 \left.\begin{matrix}\mathbf a,\alpha\\\beta,\mathbf b\end{matrix}\right|z\right).
\label{eq:euler-one-infty}
\end{align}
\end{theorem}
\begin{proof}
For \eqref{eq:euler-zero-one}, perform the calculation
\eqref{eq:beta-inner} on $[\epsilon,1-\epsilon]$, where all sums and integrals
are finite.  The beta identity itself is proved by rational quadrature and the
Gamma/Beta contract.  The endpoint majorants send $\epsilon$ to the requested
precision.  For \eqref{eq:euler-one-infty}, substitute $x=1/t$; the inner
integral becomes a beta integral whose Gamma quotient is
$\Gamma(\beta-s)/\Gamma(\alpha-s)$, inserting the bottom numerator parameter
$\beta$ and top denominator parameter $\alpha$.  The supplied Fubini budget
justifies the finite order exchange.
\end{proof}

\subsection{Convolution of two \texorpdfstring{$G$}{G}-functions}

Let the second $G$ have parameters $\mathbf c$ (length $\sigma$),
$\mathbf d$ (length $\tau$), and indices $\mu,\nu$.  The source gives two
equivalent concatenations.

\begin{theorem}[Meijer convolution]\label{thm:convolution}
Under a common Mellin strip, positive scales $\eta,\omega$, and a convolution
certificate,
\begin{align}
&\int_0^\infty
G_{p,q}^{m,n}(\eta x)G_{\sigma,\tau}^{\mu,\nu}(\omega x)\dd x\notag\\
&=\frac1\eta
G_{q+\sigma,p+\tau}^{n+\mu,m+\nu}\!\left(
 \left.\begin{matrix}-b_1,\ldots,-b_m,\mathbf c,-b_{m+1},\ldots,-b_q\\
 -a_1,\ldots,-a_n,\mathbf d,-a_{n+1},\ldots,-a_p
 \end{matrix}\right|\frac\omega\eta\right)\label{eq:conv-one}\\
&=\frac1\omega
G_{p+\tau,q+\sigma}^{m+\nu,n+\mu}\!\left(
 \left.\begin{matrix}a_1,\ldots,a_n,-\mathbf d,a_{n+1},\ldots,a_p\\
 b_1,\ldots,b_m,-\mathbf c,b_{m+1},\ldots,b_q
 \end{matrix}\right|\frac\eta\omega\right).
\label{eq:conv-two}
\end{align}
\end{theorem}
\begin{proof}
Insert a finite Mellin--Barnes approximation for the first factor.  The inner
$x$-integral is the Mellin transform of the second factor at $1-s$, supplied by
\cref{thm:mellin}.  Its Gamma factors concatenate with those already present;
reading the four blocks gives \eqref{eq:conv-one}.  Insert the second factor
instead to obtain \eqref{eq:conv-two}; alternatively apply inversion and the
scale change $x\mapsto(\eta/\omega)x$.  The convolution certificate bounds the
four discarded regions in the finite $(x,t)$ rectangle, so the order exchange
is an iPRA finite-sum equality plus error inequalities.
\end{proof}

\subsection{Products with rational powers}

The introductory source claim allows arguments $cx^\gamma$ and $dx^\delta$.
The following finite compiler proves its mathematically precise form.

\begin{theorem}[Rational-power product compiler]\label{thm:rational-power-product}
Let $\gamma,\delta>0$ be rational, write
$\gamma/\delta=A/B$ in lowest terms with $A,B\in\N_{>0}$, and let $c,d>0$.
Under a common endpoint, branch, Mellin, and Fubini certificate, every integral
\begin{equation}\label{eq:rational-power-product}
 \int_0^\infty x^{\lambda-1}G_1(cx^\gamma)G_2(dx^\delta)\dd x
\end{equation}
is a computable constant times one ordinary Meijer $G$-function.  Its four
parameter blocks, argument, and constant are produced by a finite
primitive-recursive list transformation.  Taking $\lambda=1$ proves the
product-integral assertion on the source page.
\end{theorem}
\begin{proof}
Insert
$G_1(cx^\gamma)=(2\pi\ii)^{-1}\int\Phi_1(s)c^sx^{\gamma s}\dd s$.
The inner integral is the power-substituted Mellin transform
\[
 \int_0^\infty x^{\lambda+\gamma s-1}G_2(dx^\delta)\dd x
 =\frac1\delta d^{-(\lambda+\gamma s)/\delta}
 M_2\!\left(\frac{\lambda+\gamma s}{\delta}\right).
\]
This is first performed on a finite rectangle; the certificate controls all
discarded pieces.  Put $s=Bu$.  Gamma factors from $\Phi_1(Bu)$ have slope
$\pm B$, while those from $M_2(\lambda/\delta+Au)$ have slope $\pm A$.
Gauss multiplication replaces them by unit-slope factors, for example
\begin{align*}
 \Gamma(\alpha+Au)
 &= (2\pi)^{(1-A)/2}A^{\alpha+Au-1/2}
    \prod_{r=0}^{A-1}\Gamma\!\left(u+\frac{\alpha+r}{A}\right),\\
 \Gamma(\beta-Bu)
 &= (2\pi)^{(1-B)/2}B^{\beta-Bu-1/2}
    \prod_{r=0}^{B-1}\Gamma\!\left(\frac{\beta+r}{B}-u\right).
\end{align*}
The inverted formulas handle reciprocal Gamma factors.  Tagging each result by
slope sign and numerator/denominator position sorts the finite list into the
four standard blocks.  All $u$-dependent powers combine into one new argument;
the rest form the prefactor.  Hence the resulting integral is one ordinary
Meijer $G$-integral.  List expansion, sorting, affine parameter formation, and
exponent collection are bounded loops, and tail transport uses the same
finite Gamma-multiplication bounds as \cref{thm:multiplication}.
\end{proof}

\subsection{Laplace transforms}

The elementary inner integral is
\begin{equation}\label{eq:laplace-inner}
 \int_0^\infty \e^{-\omega x}x^{s-\alpha}\dd x
 =\omega^{\alpha-s-1}\Gamma(1-\alpha+s),
 \qquad \Re\omega>0.
\end{equation}

\begin{theorem}[Laplace and inverse Laplace]\label{thm:laplace}
Under endpoint and Fubini certificates and $\Re\omega>0$,
\begin{equation}\label{eq:laplace-main}
\int_0^\infty \e^{-\omega x}x^{-\alpha}G_{p,q}^{m,n}(\eta x)\dd x
=\omega^{\alpha-1}
G_{p+1,q}^{m,n+1}\!\left(
 \left.\begin{matrix}\alpha,\mathbf a\\\mathbf b\end{matrix}\right|
 \frac\eta\omega\right).
\end{equation}
If a Bromwich-line certificate places $\Re\omega=c$ to the right of all poles,
then the inverse formula is
\begin{equation}\label{eq:inverse-laplace}
 x^{-\alpha}G_{p,q+1}^{m,n}(\mathbf a;\mathbf b,\alpha\mid\eta x)
 =\frac1{2\pi\ii}\int_{c-\ii\infty}^{c+\ii\infty}
 \e^{\omega x}\omega^{\alpha-1}
 G_{p,q}^{m,n}(\mathbf a;\mathbf b\mid\eta/\omega)\dd\omega.
\end{equation}
\end{theorem}
\begin{proof}
Equation \eqref{eq:laplace-inner} inserts the top numerator parameter $\alpha$
and extracts $\omega^{\alpha-1}$; the remaining power changes the argument to
$\eta/\omega$.  Finite Fubini proves \eqref{eq:laplace-main}.  The inverse is
finite Bromwich/Fourier inversion after $\omega=c+\ii t$, with exactly the same
Fej\'er-kernel method as \cref{thm:mellin}.  The line and tail bounds are part
of the certificate, rather than hidden in the phrase ``as usual''.
\end{proof}

\begin{theorem}[Quadratic-argument Laplace formula]\label{thm:quadratic-laplace}
For $\Re\omega>0$ and the same type of certificates,
\begin{equation}\label{eq:quadratic-laplace}
 \int_0^\infty \e^{-\omega x}G_{p,q}^{m,n}(\eta x^2)\dd x
 =\frac1{\sqrt\pi\,\omega}
 G_{p+2,q}^{m,n+2}\!\left(
 \left.\begin{matrix}0,\frac12,\mathbf a\\\mathbf b\end{matrix}\right|
 \frac{4\eta}{\omega^2}\right).
\end{equation}
\end{theorem}
\begin{proof}
Inside a finite contour use
$\int_0^\infty\e^{-\omega x}x^{2s}\dd x=
\omega^{-2s-1}\Gamma(2s+1)$ and the duplication identity
\[
 \Gamma(2s+1)=\frac{2^{2s}}{\sqrt\pi}\Gamma(s+1)\Gamma(s+\tfrac12).
\]
The two new Gamma factors correspond to top parameters $0$ and $1/2$, and
$2^{2s}\eta^s\omega^{-2s}=(4\eta/\omega^2)^s$.
\end{proof}

\subsection{Transform kernels: what is proved}

A pair $(k,h)$ is a transform pair on a coded test class $\mathcal T$ when the
two finite-window operators
\[
 (K_Rf)(z)=\int_{1/R}^{R}k(z,y)f(y)\dd y,
 \qquad
 (H_Rg)(z)=\int_{1/R}^{R}h(z,y)g(y)\dd y
\]
have quadrature and tail certificates, and
$H_{R(k)}K_{R(k)}f$ is within $2^{-k}$ of $f$ on the requested rational compact
set.  This definition is first-order once $f,k,h,R$ are fixed PR names.

\begin{theorem}[Certified kernel principle]\label{thm:kernel-principle}
Suppose the Mellin multipliers of $k$ and $h$ multiply to $1$ on a rational
strip, and suppose the test function and kernels carry common Mellin, endpoint,
and Fubini certificates.  Then $(k,h)$ is a transform pair in the finite sense
above.
\end{theorem}
\begin{proof}
Take finite Mellin transforms.  Convolution becomes multiplication by
\cref{thm:convolution}; the two multipliers cancel algebraically.  Apply the
finite Mellin inversion theorem and add the four explicit truncation errors.
\end{proof}

The Narain kernels on the source page have reciprocal Gamma-ratio multipliers
once
$n-p=m-q>0$ and
$\sum a_j+\sum b_j=\sum c_j+\sum d_j$.  Therefore they instantiate
\cref{thm:kernel-principle} whenever the further endpoint conditions are
supplied.  If $p=q$, $m=n$, $a_j+b_j=0$, and $c_j+d_j=0$, parameter inversion
and reordering show $k=h$; the Hankel kernel is the stated special case.  The
Wimp kernels are treated identically after expanding the auxiliary $A$-function
into its $G$-kernel.  These are status \statusQ/\statusR results because the
source does not specify a maximal test class.

The generalized Laplace pair follows from \cref{thm:kernel-principle} with its
stated balance
\[
 (q+1-p)\frac{\rho}{2\gamma}=\sum_ja_j-\sum_jb_j,
 \qquad \gamma>0,\quad p\le q.
\]
At $\gamma=1/2$, $\rho=0$, and $p=q=0$, the kernel reduces by
\cref{thm:exp-case} to the ordinary Laplace kernel.  The two Meijer transforms
with $K_\nu/I_\nu$ and with Whittaker $W/M$ are direct specializations.  For the
second, set $\gamma=1/2$, $\rho=-m-k$, $p=q=1$, $a_1=m-k$, and $b_1=2m$.
The forward-kernel identification is the residue identity
\begin{equation}\label{eq:g-u-whittaker}
 G_{1,2}^{2,0}\!\left(
  \left.\begin{matrix}a\\b_1,b_2\end{matrix}\right|x\right)
 =x^{b_2}\e^{-x}U(a-b_1,1+b_2-b_1,x),
\end{equation}
with a common branch.  Taking $(a,b_1,b_2)=(m-k,0,2m)$ and using
\[
 W_{k+1/2,m}(x)=x^{m+1/2}\e^{-x/2}U(m-k,1+2m,x)
\]
gives exactly the source's $x^{-k-1/2}\e^{-x/2}W_{k+1/2,m}(x)$ kernel after
multiplying by the generalized-Laplace power $x^{-m-k}$.  For the inverse
kernel, \cref{thm:hyp-to-g} gives
\[
 G_{1,2}^{1,1}\!\left(
 \left.\begin{matrix}k-m\\0,-2m\end{matrix}\right|-x\right)
 =\frac{\Gamma(1+m-k)}{\Gamma(1+2m)}
 {}_1F_1(1+m-k;1+2m;x).
\]
Combining this with
$M_{k-1/2,m}(x)=x^{m+1/2}\e^{-x/2}
 {}_1F_1(1+m-k;1+2m;x)$ produces the displayed inverse kernel and its Gamma
normalization.  Formula \eqref{eq:g-u-whittaker} follows from the same two
residue families as the upper incomplete Gamma case, with a factorial tail.
Thus both Meijer pairs have status \statusQ under the inherited certificates.
An unrestricted theorem for every ``suitable'' function is intentionally not
asserted.

\section{Elementary functions and the MathWorld special cases}\label{sec:elementary-cases}

Every identity in this section is proved by the same finite mechanism.  Insert
the relevant parameters in the one-pole Slater formula, cancel equal
Pochhammer factors, and compare the resulting coefficient recurrence with the
standard primitive-recursive power series.  Thus no appeal to an abstract
identity theorem is needed.  Branch statements below mean equality of the
certified names on the indicated branch boxes.

\subsection{Exponential, trigonometric, and hyperbolic functions}

\begin{theorem}[Exponential and circular functions]\label{thm:exp-case}
On compatible branches,
\begin{align}
 \e^x
 &=G_{0,1}^{1,0}\!\left(\left.\begin{matrix}-\\0\end{matrix}\right|-x\right),
 \label{eq:exp-g}\\
 \cos x
 &=\sqrt\pi\,G_{0,2}^{1,0}\!\left(
       \left.\begin{matrix}-\\0,\frac12\end{matrix}\right|\frac{x^2}{4}\right),
 \label{eq:cos-g}\\
 \sin x
 &=\sqrt\pi\,G_{0,2}^{1,0}\!\left(
       \left.\begin{matrix}-\\\frac12,0\end{matrix}\right|\frac{x^2}{4}\right).
 \label{eq:sin-g}
\end{align}
For \eqref{eq:sin-g} the square-root branch is normalized by
$\sqrt{x^2}=x$, for example on $-\pi/2<\Arg x\le\pi/2$.
\end{theorem}
\begin{proof}
For \eqref{eq:exp-g}, \cref{thm:slater-right} has one pole family, no
prefactor, and hypergeometric part ${}_0F_0(-x)=\e^{-x}$ before replacing the
argument by $-x$.

For the cosine parameters, the residue at $s=k$ gives
\[
 G_{0,2}^{1,0}\!\left(
       \left.\begin{matrix}-\\0,\frac12\end{matrix}\right|z\right)
 =\frac1{\sqrt\pi}{}_0F_1\!\left(;\frac12;-z\right).
\]
The $k$th coefficient after $z=x^2/4$ is
\[
 \frac{(-1)^k x^{2k}}{4^k k!\,\Gamma(k+\frac12)}
 =\frac{(-1)^k x^{2k}}{\sqrt\pi(2k)!},
\]
where the duplication identity
$\Gamma(k+\frac12)=\sqrt\pi(2k)!/(4^k k!)$ is an exact finite factorial
identity.  Multiplication by $\sqrt\pi$ gives the cosine series.

Similarly,
\[
 G_{0,2}^{1,0}\!\left(
       \left.\begin{matrix}-\\\frac12,0\end{matrix}\right|z\right)
 =\frac{2z^{1/2}}{\sqrt\pi}{}_0F_1\!\left(;\frac32;-z\right).
\]
With $z=x^2/4$ and the stated branch, the coefficient is
$(-1)^kx^{2k+1}/[\sqrt\pi(2k+1)!]$.  The hypergeometric tail bounds from
\cref{thm:hyp-conv} are also the usual sine and cosine tail bounds, so the
coefficient equalities lift directly to ball equality at every precision.
\end{proof}

\begin{corollary}[Hyperbolic functions]\label{cor:hyperbolic-g}
On the branches displayed on the source page,
\begin{align}
 \cosh x&=\sqrt\pi\,G_{0,2}^{1,0}\!\left(
  \left.\begin{matrix}-\\0,\frac12\end{matrix}\right|-\frac{x^2}{4}\right),
 \label{eq:cosh-g}\\
 \sinh x&=-\ii\sqrt\pi\,G_{0,2}^{1,0}\!\left(
  \left.\begin{matrix}-\\\frac12,0\end{matrix}\right|-\frac{x^2}{4}\right).
 \label{eq:sinh-g}
\end{align}
\end{corollary}
\begin{proof}
Replace $x$ by $\ii x$ in \eqref{eq:cos-g}--\eqref{eq:sin-g}.  This is not a
formal appeal to complex substitution: term by term, $(-1)^k(\ii x)^{2k}=x^{2k}$
and $(-1)^k(\ii x)^{2k+1}=\ii x^{2k+1}$.  The branch condition in
\eqref{eq:sinh-g} is precisely the certificate that
$\sqrt{-x^2}=\ii x$.
\end{proof}

\subsection{Inverse circular functions and the logarithm}

\begin{theorem}[Arcsine, arctangent, and arccotangent]\label{thm:inverse-circular}
With the indicated branch normalizations,
\begin{align}
 \arcsin x
 &=\frac{-\ii}{2\sqrt\pi}G_{2,2}^{1,2}\!\left(
   \left.\begin{matrix}1,1\\\frac12,0\end{matrix}\right|-x^2\right),
 \label{eq:arcsin-g}\\
 \arctan x
 &=\frac12G_{2,2}^{1,2}\!\left(
   \left.\begin{matrix}\frac12,1\\\frac12,0\end{matrix}\right|x^2\right),
 \label{eq:arctan-g}\\
 \operatorname{arccot}x
 &=\frac12G_{2,2}^{2,1}\!\left(
   \left.\begin{matrix}\frac12,1\\\frac12,0\end{matrix}\right|x^2\right).
 \label{eq:arccot-g}
\end{align}
Here $\operatorname{arccot}x$ is the branch selected by
$\operatorname{arccot}x=\arctan(1/x)$, not an unspecified global convention.
\end{theorem}
\begin{proof}
The right-residue formula and elementary Gamma cancellation give
\begin{align*}
G_{2,2}^{1,2}\!\left(
 \left.\begin{matrix}1,1\\\frac12,0\end{matrix}\right|z\right)
 &=2\sqrt\pi\,z^{1/2}
 {}_2F_1\!\left(\frac12,\frac12;\frac32;-z\right),\\
G_{2,2}^{1,2}\!\left(
 \left.\begin{matrix}\frac12,1\\\frac12,0\end{matrix}\right|z\right)
 &=2z^{1/2}
 {}_2F_1\!\left(1,\frac12;\frac32;-z\right).
\end{align*}
On the first branch $(-x^2)^{1/2}=\ii x$; hence the first expression becomes
$x\,{}_2F_1(1/2,1/2;3/2;x^2)$ after multiplying by
$-\ii/(2\sqrt\pi)$.  This is the coefficient definition of $\arcsin x$.
The second becomes
$x\,{}_2F_1(1,1/2;3/2;-x^2)$, whose coefficient recurrence is
$(-1)^kx^{2k+1}/(2k+1)$.

For \eqref{eq:arccot-g}, apply the inversion theorem
\eqref{eq:inversion}: the displayed $G_{2,2}^{2,1}(x^2)$ equals the
$G_{2,2}^{1,2}$ in \eqref{eq:arctan-g} at $x^{-2}$.  Thus its half is
$\arctan(1/x)$ on the chosen branch.  All three statements extend from their
convergent local series exactly as certified names, using either the finite
contour presentation or a chain of overlapping rational discs with explicit
Taylor remainders.
\end{proof}

\begin{theorem}[Logarithm and rational function]\label{thm:log-rational}
For a branch box avoiding $(-\infty,-1]$ and the relevant poles,
\begin{align}
 G_{2,2}^{1,2}\!\left(
  \left.\begin{matrix}1,1\\1,0\end{matrix}\right|z\right)
 &=\Log(1+z),
 \label{eq:log-g}\\
 G_{2,2}^{1,2}\!\left(
  \left.\begin{matrix}1,1\\1,1\end{matrix}\right|z\right)
 &=\frac{z}{1+z}.
 \label{eq:rational-g}
\end{align}
These prove MathWorld formulas (4) and (5), with the missing domain
qualification restored.
\end{theorem}
\begin{proof}
Slater's formula gives, respectively,
\[
 z\,{}_2F_1(1,1;2;-z),
 \qquad z\,{}_2F_1(1,1;1;-z).
\]
The first has coefficient $(-1)^kz^{k+1}/(k+1)$ and therefore is
$\Log(1+z)$.  In the second, the common parameter $1$ cancels and
${}_1F_0(1;;-z)=(1+z)^{-1}$.  The tail estimates are geometric on
$\abs z<1$; outside that disc, the certified contour or the algebraic identity
$(1+z)F=z$ supplies continuation without excluded-middle reasoning.
\end{proof}

\subsection{The four MathWorld formulas}

\begin{theorem}[MathWorld special cases]\label{thm:mathworld-specials}
In addition to \eqref{eq:log-g}--\eqref{eq:rational-g},
\begin{align}
G_{0,2}^{1,0}\!\left(
 \left.\begin{matrix}-\\0,\frac12\end{matrix}\right|\frac z2\right)
 &=\frac{\cos\sqrt{2z}}{\sqrt\pi},
 \label{eq:mathworld-cos}\\
G_{1,0}^{0,1}\!\left(
 \left.\begin{matrix}1-a\\-\end{matrix}\right|z\right)
 &=z^{-a}\e^{-1/z},
 \label{eq:mathworld-recip-exp}\\
G_{0,2}^{1,0}\!\left(
 \left.\begin{matrix}-\\0,\frac12\end{matrix}\right|\frac z2;\frac12\right)
 &=\frac{\cos z}{\sqrt\pi}.
 \label{eq:mathworld-r-cos}
\end{align}
The last line uses the generalized fourth parameter of
\eqref{eq:r-general}.
\end{theorem}
\begin{proof}
Equation \eqref{eq:mathworld-cos} is \eqref{eq:cos-g} with
$x=\sqrt{2z}$.  The left-residue formula for the single top pole in
\eqref{eq:mathworld-recip-exp} is
$z^{-a}{}_0F_0(-1/z)$.

For \eqref{eq:mathworld-r-cos}, take residues in the generalized integral.
The $k$th term is
\[
 \frac{(-1)^k(z/2)^{2k}}{k!\,\Gamma(k+\frac12)}
 =\frac{(-1)^kz^{2k}}{\sqrt\pi(2k)!}.
\]
The equality is therefore the cosine series.  This also explains why the
fourth parameter cannot be discarded without changing the branch and scale.
\end{proof}

\subsection{Step functions and a source correction}

\begin{theorem}[Positive-real step kernels]\label{thm:step-g}
For positive real $x$ with $x\ne1$,
\begin{align}
G_{1,1}^{1,0}\!\left(\left.\begin{matrix}1\\0\end{matrix}\right|x\right)
 &=\begin{cases}1,&0<x<1,\\0,&x>1,\end{cases}
 \label{eq:step-in}\\
G_{1,1}^{0,1}\!\left(\left.\begin{matrix}1\\0\end{matrix}\right|x\right)
 &=\begin{cases}0,&0<x<1,\\1,&x>1.\end{cases}
 \label{eq:step-out}
\end{align}
At $x=1$ the value depends on the chosen Heaviside convention.  For real
signed input the source's formulas for $H(1-|x|)$ and $H(|x|-1)$ must be read
with $G(|x|)$, or restricted to $x>0$; the literal ``for all $x$'' statement
with $G(x)$ is not branch-independent.
\end{theorem}
\begin{proof}
For $0<x<1$, the right-loop residue formula for $G_{1,1}^{1,0}$ has one pole
and reduces to ${}_1F_0(0;;-x)=1$, whereas $G_{1,1}^{0,1}$ has no enclosed
right poles and is zero.  For $x>1$, the left-loop roles reverse: the first has
no left pole and the second has one residue equal to $1$.  The closing-side
bounds are geometric, so this is a finite remainder proof at every requested
precision.  No claim is made at the discontinuity without a separate coded
boundary convention.
\end{proof}

\subsection{The six-term representation of the identity function}

The source also prints a six-term expression involving
$G_{1,3}^{3,0}$ and $G_{2,4}^{2,0}$ which equals $\abs x$, followed by a factor
$\operatorname{sgn}x$.  Each summand is a confluent case, so expanding it by
naive nonconfluent Slater formulas is invalid.  Applying the Laurent-jet
algorithm of \cref{def:jet}, then the contiguous relations of
\cref{sec:derivatives}, reduces the bracket to the unique solution $F(t)$ of
\[
 2tF'(t)-F(t)=0,\qquad F(1)=1,
\]
namely $F(t)=\sqrt t$ on $t>0$.  The reduction is a finite rational certificate
in the parameter lists; the companion numerical script independently checks
the printed combination at several positive points.

\begin{proposition}[Corrected scope of the six-term identity]
For real $x\ne0$, with every $G$ evaluated on the positive argument $x^2$ and
with the real positive square root, the six-term bracket printed by the source
is $\abs x$, hence multiplication by $\operatorname{sgn}x$ gives $x$.  At
$x=0$ the statement means the common continuous limit; the individual printed
terms need not be defined there.
\end{proposition}

This proposition is assigned status \statusQ rather than \statusP in the
ledger: the report specifies the exact differential/contiguous certificate and
checks it computationally, but does not print the several-page normalized
certificate term by term.  It is therefore not used as a dependency of any
other result.

\section{Incomplete gamma, Bessel, and Lerch functions}\label{sec:higher-cases}

The web pages state these formulas with ``for all $x$''.  Constructively, each
formula has three pieces of domain data: a branch for the powers of $x$, an
apartness certificate for parameter denominators, and either a convergent
series disc or a contour/tail certificate.  The formulas below include those
qualifications explicitly.

\subsection{Lower and upper incomplete gamma functions}

For $\Re\alpha>0$ and a path from $0$ to $x$ contained in a certified branch
sector, define
\[
 \gamma(\alpha,x)=\int_0^x t^{\alpha-1}\e^{-t}\dd t,
 \qquad
 \Gamma(\alpha,x)=\int_x^\infty t^{\alpha-1}\e^{-t}\dd t.
\]
The integrals mean finite rational quadrature with endpoint and tail bounds.
For other admissible parameters they are continued by the recurrence
$\gamma(\alpha+1,x)=\alpha\gamma(\alpha,x)-x^\alpha\e^{-x}$ and the analogous
upper relation.

\begin{theorem}[Incomplete gamma representations]\label{thm:incomplete-gamma-g}
On every such certified domain,
\begin{align}
 \gamma(\alpha,x)
 &=G_{1,2}^{1,1}\!\left(
   \left.\begin{matrix}1\\\alpha,0\end{matrix}\right|x\right),
 \label{eq:lower-gamma-g}\\
 \Gamma(\alpha,x)
 &=G_{1,2}^{2,0}\!\left(
   \left.\begin{matrix}1\\\alpha,0\end{matrix}\right|x\right).
 \label{eq:upper-gamma-g}
\end{align}
\end{theorem}
\begin{proof}
The one-pole Slater formula gives
\begin{equation}\label{eq:lower-hyp}
G_{1,2}^{1,1}\!\left(
 \left.\begin{matrix}1\\\alpha,0\end{matrix}\right|x\right)
 =\frac{x^\alpha}{\alpha}{}_1F_1(\alpha;\alpha+1;-x).
\end{equation}
Its $k$th term is $(-1)^kx^{\alpha+k}/[k!(\alpha+k)]$.  Integrating the
exponential series term by term on a finite segment gives the same series for
$\gamma(\alpha,x)$, with a factorial tail bound, proving
\eqref{eq:lower-gamma-g}.

For the upper function, assume first that $\alpha\notin\Z$.  The two right pole
families give
\[
G_{1,2}^{2,0}\!\left(
 \left.\begin{matrix}1\\\alpha,0\end{matrix}\right|x\right)
 =\Gamma(\alpha)-\frac{x^\alpha}{\alpha}
 {}_1F_1(\alpha;\alpha+1;-x).
\]
Indeed the $b=0$ residue is the constant $\Gamma(\alpha)$, while
$\Gamma(-\alpha)/\Gamma(1-\alpha)=-1/\alpha$ for the $b=\alpha$ family.
The right side is $\Gamma(\alpha)-\gamma(\alpha,x)=\Gamma(\alpha,x)$.
When the pole families coalesce, the finite Laurent-jet algorithm computes the
same cancellation directly; alternatively both sides satisfy
$\partial_xF=-x^{\alpha-1}\e^{-x}$ and have the same certified tail value at
infinity.  This proves the integer and resonant cases without an illicit
nonconfluent formula.
\end{proof}

\subsection{Bessel \texorpdfstring{$J,Y,I,K$}{J, Y, I, K}}

\begin{theorem}[Bessel representations]\label{thm:bessel-g}
With $\sqrt{x^2}=x$ on the indicated branch,
\begin{align}
 J_\nu(x)
 &=G_{0,2}^{1,0}\!\left(
  \left.\begin{matrix}-\\\frac\nu2,-\frac\nu2\end{matrix}\right|
  \frac{x^2}{4}\right),
 \label{eq:J-g}\\
 Y_\nu(x)
 &=G_{1,3}^{2,0}\!\left(
  \left.\begin{matrix}-\frac{\nu+1}{2}\\
  \frac\nu2,-\frac\nu2,-\frac{\nu+1}{2}\end{matrix}\right|
  \frac{x^2}{4}\right),
 \label{eq:Y-g}\\
 I_\nu(x)
 &=\ii^{-\nu}G_{0,2}^{1,0}\!\left(
  \left.\begin{matrix}-\\\frac\nu2,-\frac\nu2\end{matrix}\right|
  -\frac{x^2}{4}\right),
 \label{eq:I-g}\\
 K_\nu(x)
 &=\frac12G_{0,2}^{2,0}\!\left(
  \left.\begin{matrix}-\\\frac\nu2,-\frac\nu2\end{matrix}\right|
  \frac{x^2}{4}\right).
 \label{eq:K-g}
\end{align}
For $J,Y,K$ one may take $-\pi/2<\Arg x\le\pi/2$; for the displayed $I$
formula one may take $-\pi<\Arg x\le0$.  Other sectors follow by certified
branch transport.
\end{theorem}
\begin{proof}
For \eqref{eq:J-g}, Slater's formula is
\[
 \frac{(x/2)^\nu}{\Gamma(\nu+1)}
 {}_0F_1\!\left(;\nu+1;-\frac{x^2}{4}\right),
\]
which is the defining Bessel series.  Replacing $x$ by $\ii x$ and tracking the
power branch gives \eqref{eq:I-g}.

Suppose first that $\nu\notin\Z$.  The two residues for
$G_{0,2}^{2,0}(z\mid\nu/2,-\nu/2)$ give
\[
 \Gamma(-\nu)z^{\nu/2}{}_0F_1(;1+\nu;z)
 +\Gamma(\nu)z^{-\nu/2}{}_0F_1(;1-\nu;z).
\]
At $z=x^2/4$, one half of this is the standard constructive definition of
$K_\nu(x)$.  Equivalently it follows from
$K_\nu=\pi(I_{-\nu}-I_\nu)/(2\sin\pi\nu)$ and the reflection formula.  At
integer $\nu$ the two pole families coalesce; \cref{def:jet} returns the usual
power--log series and proves the limit without choosing a sequence of
nonintegers.

For \eqref{eq:Y-g}, the two nonconfluent residues reduce to
\[
 \cot(\pi\nu)J_\nu(x)-\csc(\pi\nu)J_{-\nu}(x)
 =Y_\nu(x).
\]
Here the first coefficient follows from
$\Gamma(\nu+\tfrac32)\Gamma(-\nu-\tfrac12)=-\pi/\cos\pi\nu$ and
$\Gamma(-\nu)\Gamma(1+\nu)=-\pi/\sin\pi\nu$; the second follows from
$\Gamma(\nu)\Gamma(1-\nu)=\pi/\sin\pi\nu$.  Integer orders again use the
Laurent-jet normal form.  These calculations are finite Gamma identities plus
hypergeometric coefficient recurrences and hence fit the iPRA contracts.
\end{proof}

\subsection{The Lerch transcendent at integer order}

For a fixed $n\in\N$, $\abs x<1$, and $a$ quantitatively apart from
$0,-1,-2,\ldots$, put
\[
 \Phi(x,n,a)=\sum_{k=0}^\infty\frac{x^k}{(a+k)^n},
 \qquad
 \Phi(x,-n,a)=\sum_{k=0}^\infty(a+k)^n x^k.
\]
The geometric tail, multiplied by a fixed polynomial in $k$, gives an explicit
primitive-recursive modulus.

\begin{theorem}[Lerch representations]\label{thm:lerch-g}
For $n\in\N$ and initially $\abs x<1$,
\begin{align}
 \Phi(x,n,a)
 &=G_{n+1,n+1}^{1,n+1}\!\left(
 \left.\begin{matrix}0,\underbrace{1-a,\ldots,1-a}_{n}\\
 0,\underbrace{-a,\ldots,-a}_{n}\end{matrix}\right|-x\right),
 \label{eq:lerch-positive}\\
 \Phi(x,-n,a)
 &=G_{n+1,n+1}^{1,n+1}\!\left(
 \left.\begin{matrix}0,\underbrace{-a,\ldots,-a}_{n}\\
 0,\underbrace{1-a,\ldots,1-a}_{n}\end{matrix}\right|-x\right).
 \label{eq:lerch-negative}
\end{align}
Outside the unit disc these formulas specify a branch continuation under the
usual $G$-function certificates.  The source's phrase ``for all $x$'' is not a
single-valued, branch-free assertion.
\end{theorem}
\begin{proof}
The single right residue in \eqref{eq:lerch-positive} yields
\[
 a^{-n}{}_{n+1}F_n\!\left(
 \begin{matrix}1,a,\ldots,a\\a+1,\ldots,a+1\end{matrix};x\right).
\]
Since $(a)_k/(a+1)_k=a/(a+k)$ and $(1)_k=k!$, its $k$th coefficient is
$x^k/(a+k)^n$.  For \eqref{eq:lerch-negative}, the prefactor is $a^n$ and the
hypergeometric quotient is
$[(a+1)_k/(a)_k]^n=[(a+k)/a]^n$, giving $(a+k)^nx^k$.  Repeated parameters do
not create repeated right poles because $m=1$; the repeated top parameters are
ordinary finite products.  Thus the proof is a direct coefficient identity
with a geometric-polynomial tail.
\end{proof}

\subsection{Parameter derivatives of incomplete gamma functions}

The final sentence of the Wikipedia list says that derivatives with respect to
$\alpha$ can also be written with $G$-functions.  A useful constructive form is
more explicit than that sentence.

\begin{theorem}[All fixed parameter derivatives]\label{thm:gamma-parameter-derivatives}
Let $r\in\N$ be externally fixed, $\Re\alpha\ge\epsilon>0$, and let $x$ lie in
a certified logarithm branch.  Then
\begin{align}
 \frac{\partial^r}{\partial\alpha^r}\gamma(\alpha,x)
 &=x^\alpha\sum_{\ell=0}^{r}
 \binom{r}{\ell}(\Log x)^{r-\ell}(-1)^\ell\ell!\,
 \alpha^{-\ell-1}
 {}_{\ell+1}F_{\ell+1}\!\left(
 \begin{matrix}\alpha,\ldots,\alpha\\
 \alpha+1,\ldots,\alpha+1\end{matrix};-x\right).
 \label{eq:gamma-param-derivative}
\end{align}
There are $\ell+1$ copies in each parameter row.  Every hypergeometric term in
\eqref{eq:gamma-param-derivative} has a Meijer $G$ representation by
\cref{thm:hyp-to-g}; $\Log x$ has \eqref{eq:log-g}.  The corresponding upper
incomplete-gamma derivative is
\[
 \partial_\alpha^r\Gamma(\alpha,x)
 =\Gamma^{(r)}(\alpha)-\partial_\alpha^r\gamma(\alpha,x),
\]
with $\Gamma^{(r)}$ supplied by the fixed-order Gamma contract.  Equivalently,
a single confluent $G$ expression can be produced by the Laurent-jet compiler.
\end{theorem}
\begin{proof}
Differentiate the uniformly bounded series
\[
 \gamma(\alpha,x)=x^\alpha\sum_{k=0}^\infty
 \frac{(-x)^k}{k!(\alpha+k)}.
\]
The finite Leibniz rule gives
\[
 \frac{\partial^r}{\partial\alpha^r}
 \frac{x^\alpha}{\alpha+k}
 =x^\alpha\sum_{\ell=0}^{r}\binom{r}{\ell}
 (\Log x)^{r-\ell}\frac{(-1)^\ell\ell!}{(\alpha+k)^{\ell+1}}.
\]
Finally,
\[
 \sum_{k=0}^\infty\frac{(-x)^k}{k!(\alpha+k)^d}
 =\alpha^{-d}{}_dF_d(\alpha,\ldots,\alpha;
                    \alpha+1,\ldots,\alpha+1;-x).
\]
For a requested precision, a common factorial tail bound controls all
$0\le\ell\le r$, so differentiation is replaced by a finite coefficient
calculation plus an explicit remainder estimate.  This is exactly the fixed
finite-use form permitted in iPRA.
\end{proof}

\subsection{What ``most special functions'' means constructively}

The hypergeometric conversion theorem already covers every fixed generalized
hypergeometric function away from polynomial singularities, and
\cref{thm:polynomial-repair} covers the polynomial cases.  Products, Mellin
convolutions, Laplace transforms, parameter shifts, derivatives, and rational
power substitutions are closed by \cref{sec:algebra,sec:derivatives,sec:transforms}.
Consequently any particular special-function identity obtained from a finite
chain of those operations compiles to a finite iPRA certificate.  This is the
precise content behind the informal claim that the $G$-function contains most
classical special functions.  It is not a quantification over an undefined
class of ``all special functions''.

\section{Lean formalization outline and path to object iPRA}\label{sec:formalization}

This section specifies the representations, dependency order, induction
invariants, and final theorem types needed for a complete formalization.  The
files in \code{lean/} mirror the plan.  They are an integration scaffold, not a
claim that the analytic development has already been checked by Lean.

\subsection{The endpoint that counts}

At the pinned ProveIt revision, the ordinary iPRA syntax is in
\code{LeanProofs.ConciseIPRA}.  Primitive-recursive symbols have type
\code{PR k}; terms have type \code{PTm n}; the axioms are \code{PRAAx}; and the
convenient induction rule is \code{PRA.qfInd}.  A completed theorem must have
the shape
\begin{lstlisting}[language={}]
import ConciseIPRA.PR
import ConciseIPRA.NativeBiInterpretation

namespace LeanProofs.ConciseIPRA

-- phi : Fm PTm 0
example : [] |- [PRAAx] phi := by
  ... -- an explicit object-language derivation

end LeanProofs.ConciseIPRA
\end{lstlisting}
After this theorem exists, \code{pra\_iff\_native} transfers it to the checked
native scalar calculus.  A theorem merely about Lean's \code{Nat},
\code{Rat}, or a host complex type is useful evidence but is not a proof in
iPRA.

Every formalized result should therefore have three visibly separate layers:
\begin{enumerate}
\item an executable host definition and regression theorem;
\item a primitive-recursive graph formula or bounded certificate checker;
\item a theorem of type \code{[] |- [PRAAx] phi} proving that checker sound.
\end{enumerate}

\subsection{Compilation rather than a universal evaluator}

The user-facing notation permits variable lists of parameters, but iPRA need
not quantify over an all-purpose program interpreter.  A small external
compiler takes a concrete theorem family---fixed $p,q,m,n$, fixed formulas for
the parameter names, and fixed branch/tail algorithms---and emits ordinary
\code{PR} symbols.  List coding, bounded loops, rational arithmetic, and
finite-polynomial operations are themselves primitive recursive.  Soundness
is proved once for their checkers.  The resulting theorem is a schema over
fixed compiled names, exactly as in the existing Gamma and Barnes-$G$ packets.

For a genuinely uniform family such as ``all rational parameter lists of
length at most $r$'', $r$ is external and the generated checker uses bounded
loops through $r$.  This avoids both higher types and a universal
primitive-recursive evaluator.

\subsection{Data types and numerical predicates}

The host scaffold uses the following records.
\begin{lstlisting}[language={}]
structure QComplex where
  reNum : Int
  reDen : Nat       -- positive
  imNum : Int
  imDen : Nat       -- positive

structure MeijerSpec where
  aNum : List QComplex
  aDen : List QComplex
  bNum : List QComplex
  bDen : List QComplex

structure BranchCert where
  centerArg : Rat
  radiusArg : Rat
  awayZero  : Rat
  cutMargin : Rat

structure EvalCert where
  truncation : Nat
  guardBits  : Nat
  tailNum    : Nat
  tailPow    : Nat
  clearance  : Rat
\end{lstlisting}
The names \code{aNum}, \code{aDen}, \code{bNum}, and \code{bDen} denote the
four parameter blocks $A^-,A^+,B^-,B^+$.  In the object language these records
are natural-number tuples obtained by pairing.

A complex ball uses the $\ell^1$ norm because it avoids square roots:
\[
 B(c,p)=\{w:\abs{\Re(w-c)}+\abs{\Im(w-c)}\le2^{-p}\}.
\]
The object predicate \code{BallEq(x,y,p)} is a cleared integer inequality
between the centers and error numerators returned by two fixed evaluators.
There is no set-valued object in the formal language.  Division is never a
primitive theorem conclusion: a reciprocal operation requires
\code{Apart(q,d)}, and all final inequalities are cross-multiplied by the
positive denominator $d$.

The basic reusable object lemmas are:
\begin{enumerate}
\item signed-integer and rational normalization, order, and absolute value;
\item complex addition, multiplication, integer powers, and finite products;
\item ball addition, multiplication, Horner evaluation, and telescoping sums;
\item reciprocal with a quantitative apartness witness;
\item dyadic exponential, logarithm, trigonometric, and square-root names;
\item bounded sums, products, search, pairing, and simultaneous recursion;
\item finite Laurent polynomials with truncated convolution.
\end{enumerate}
The checked \code{IPRADerivations} arithmetic library supplies much of the
first and sixth items.  The Gamma packet is the intended source of the
certified Gamma, reciprocal-Gamma, logarithm, and Stirling interfaces, but its
analytic contracts still have to be instantiated as object derivations rather
than imported as axioms.

\subsection{Module graph}

A dependency order minimizing circularity is
\begin{verbatim}
MeijerG/Codes
  -> MeijerG/RationalBalls
  -> MeijerG/FiniteLists
  -> MeijerG/HypergeometricTerms
  -> MeijerG/HypergeometricTail
  -> MeijerG/GammaRatio
  -> MeijerG/SlaterRight
  -> MeijerG/SlaterLeft
  -> MeijerG/LaurentJets
  -> MeijerG/ResidueEvaluator
  -> MeijerG/ContourEvaluator
  -> MeijerG/ContourResidueAgreement
  -> MeijerG/DifferentialEquation
  -> MeijerG/ParameterAlgebra
  -> MeijerG/Derivatives
  -> MeijerG/Mellin
  -> MeijerG/RationalPowerProducts
  -> MeijerG/Laplace
  -> MeijerG/TransformKernels
  -> MeijerG/Whittaker
  -> MeijerG/SpecialCases
  -> MeijerG/ObjectTargets
  -> MeijerG/Checks/Axioms
\end{verbatim}
The hypergeometric and residue modules should be completed before the contour
modules.  This makes most source-page identities available without waiting for
a general complex-integration library.

\subsection{Stage 1: exact hypergeometric terms}

For externally fixed numerator parameters $\alpha_1,\ldots,\alpha_r$ and
denominator parameters $\beta_1,\ldots,\beta_s$, define a term graph by
\begin{align*}
 T_0&=1,\\
 T_{k+1}\,(k+1)\prod_{j=1}^{s}(\beta_j+k)
 &=T_k z\prod_{i=1}^{r}(\alpha_i+k).
\end{align*}
The object implementation stores a rational complex center and a common
positive denominator.  The primary theorem is the cleared recurrence; a
quotient presentation is derived only under denominator apartness.

The partial-sum loop maintains the quantifier-free invariant
\begin{equation}\label{eq:lean-hyp-invariant}
 S_k=\sum_{j<k}T_j,
 \qquad
 D_kT_{k+1}=N_kT_k,
 \qquad
 \operatorname{Err}_k\le E_k,
\end{equation}
where all equalities are represented by integer cross products.  If a bounded
checker proves $\abs{T_{j+1}}\le\rho\abs{T_j}$ for $j\ge N$ and
$\rho\le1-2^{-h}$, the geometric tail theorem yields
\[
 \abs{\sum_{j\ge N}T_j}\le \frac{\abs{T_N}}{1-\rho}.
\]
The proof is qf induction on the finite geometric sum, followed by a cleared
inequality.  This module proves \cref{thm:hyp-conv} for every compiled series.

\subsection{Stage 2: Gamma ratios and Slater coefficients}

A Slater term contains products such as
\[
 C_h=
 \frac{\prod_{j\le m,j\ne h}\Gamma(b_j-b_h)
       \prod_{j\le n}\Gamma(1+b_h-a_j)}
      {\prod_{j>m}\Gamma(1+b_h-b_j)
       \prod_{j>n}\Gamma(a_j-b_h)}.
\]
At requested precision $p$, the compiler allocates guard bits backwards
through the finite product tree, requests Gamma balls, checks denominator
apartness, and outputs a ball for $C_h$.  The soundness invariant records a
product center, an accumulated numerator error, and a positive lower bound for
every denominator factor.

The right-Slater target is schematically
\begin{lstlisting}[language={}]
RightAdmissible spec z p ->
BallEq
  (meijerRight spec z (p + guard spec p))
  (sumOverRightPoles spec z (p + guard spec p))
  p
\end{lstlisting}
and similarly for the left expansion.  The proof has two parts:
\begin{enumerate}
\item exact residue algebra for a finite closing rectangle;
\item the explicit closing-side and series-tail inequalities from
      \cref{sec:convergence}.
\end{enumerate}
For most special cases only the second line---the already normalized
hypergeometric sum---is needed as the definition, so an early formalization can
prove all coefficient identities before contour agreement.

\subsection{Stage 3: confluent poles as finite Laurent jets}

Encode a truncated Laurent polynomial by a list indexed from $-r$ to $d$.
Provide PR functions for addition, shift, truncated convolution, and coefficient
selection.  For a pole of multiplicity $r$, the invariant after multiplying
$i$ factors is
\[
 F_i(u)=\sum_{j=-r}^{r-1}c_{i,j}u^j+u^rR_i(u),
 \qquad \abs{R_i(u)}\le M_i
 \quad(\abs u\le\rho).
\]
The update for $M_i$ is a finite Cauchy-product bound.  Multiply by the finite
jet
\[
 z^{s_0+u}=z^{s_0}\sum_{j<r}\frac{(u\Log z)^j}{j!}+u^rR_z(u)
\]
and select coefficient $u^{-1}$.  The output is a polynomial in $\Log z$ of
degree at most $r-1$, exactly as in \cref{def:jet}.

The key object theorem is not ``take a limit as parameters collide''.  It is a
finite coefficient theorem saying the emitted jet is the residue of the
certified local product.  This handles integer-order $K_\nu$, $Y_\nu$, the
six-term identity in \cref{sec:elementary-cases}, and every polynomial repair.

\subsection{Stage 4: the contour evaluator}

A contour is a finite list of rational complex vertices plus, for an unbounded
vertical contour, a truncation function $T(p)$.  On each segment, use a
composite Taylor or trapezoidal rule with a supplied derivative bound.  The
object predicate \code{ContourSafe} checks:
\begin{enumerate}
\item positive distance from every relevant pole and from zero;
\item a branch box containing the image of $z^s$;
\item the vertical estimate \eqref{eq:master-vertical-bound};
\item a quadrature bound for every segment;
\item a closing-side bound for the finite residue rectangle.
\end{enumerate}
The finite residue theorem from the ProveIt Cauchy--residue packet can then be
specialized to the Gamma-ratio integrand.  Its output is the
\code{ContourResidueAgreement} theorem.  This theorem is optional for the
coefficient identities but essential for the literal Mellin--Barnes definition
and transform formulas.

\subsection{Stage 5: differential equation without differentiation axioms}

There are two complementary formalizations.

First, for a residue series $F(z)=z^\lambda\sum c_kz^k$, calculate the
coefficient of $z^{\lambda+k}$ in
\[
 (-1)^{p-m-n}z\prod_j(\thetaop-a_j+1)F
 -\prod_j(\thetaop-b_j)F.
\]
It is exactly the cleared hypergeometric term recurrence.  A qf induction proves
that every partial polynomial satisfies the equation up to an explicit tail;
the tail tends below $2^{-p}$ by the same ratio certificate.

Second, on a finite contour, applying $\thetaop$ multiplies the kernel by $s$.
Gamma recurrence converts the operator multiplier to zero identically.  The
finite quadrature error is propagated through the fixed polynomial in $s$.
This proves the ODE on contour domains and agrees with the series theorem on
overlap.

A representative object target is
\begin{lstlisting}[language={}]
Admissible spec z p ->
BallEq
  (meijerODELeft spec z (p + guard))
  (meijerODERight spec z (p + guard))
  p
\end{lstlisting}
where both sides are finite combinations of derivative-name evaluators; no
quantification over an arbitrary differentiable function occurs.

\subsection{Stage 6: parameter algebra and derivative rules}

Cancellation, power shift, block permutations, inversion, and integer block
moves are finite identities of Gamma products.  Represent a parameter list as
a bounded array and normalize it to a multiset of tagged affine expressions.
The certificate consists of a permutation, a cancellation matching, and a
small exponent adjustment.  Its checker proves equality of the kernels by
bounded product commutativity.

For the multiplication theorem, Gauss multiplication expands each Gamma
factor into $k$ factors.  The certificate records the affine map
$a\mapsto(a+r)/k$ or $b\mapsto(b+r)/k$, the extracted power of $k$, and the
argument scale.  All exponents are rational and are checked after clearing the
common denominator.

Derivative identities use a finite polynomial in $s$ and one Gamma recurrence.
The primary formal statement is a coefficient or contour ball equality, not an
unbounded limit definition of derivative.  A generic fixed-order Taylor
contract then identifies that name with the derivative on a certified disc.

\subsection{Stage 7: Mellin, Euler, rational powers, Laplace, and transform kernels}

A finite Fubini certificate contains four truncation bounds, a common
quadrature mesh, and a rational majorant for the discarded rectangles.  The
proof first rearranges a genuinely finite double sum.  It then adds:
\begin{enumerate}
\item $x$-endpoint tails;
\item contour-frequency tails;
\item quadrature errors;
\item the inner Beta or Gamma integral error.
\end{enumerate}
The Mellin transform theorem uses the triangular Fourier/Fej\'er kernel from
\cref{def:mellin-cert}; its normalization and positivity are separate finite
lemmas.  The final theorem is parameterized by a fixed function name and its
explicit modulus.  No theorem quantifies over ``all integrable functions''.

The Laplace theorem is easier: the inner integral is a Gamma contract and the
remaining Gamma product is normalized syntactically to the target $G$.
Convolution concatenates two parameter arrays after applying the Mellin theorem
at $1-s$.  These normalized-list equalities should be checked by an executable
certificate, avoiding large manual rewrites in Lean.

For rational powers $x^\gamma,x^\delta$, normalize
$\gamma/\delta=A/B$, substitute $s=Bu$, and expand every slope-$A$ or
slope-$B$ Gamma factor with the finite Gauss multiplication theorem.  The
certificate stores the expanded parameter blocks, the collected rational
power, and the prefactor.  Its soundness theorem is a bounded-list induction
plus the existing Gamma-multiplication contracts.

Transform-pair modules never quantify over an unspecified function space.  A
fixed test-function code carries endpoint, Mellin, Fubini, and inversion
moduli; the checker verifies that the two finite Mellin multipliers multiply to
one.  The Whittaker module separately compiles the $U$ and ${}_1F_1$
reductions used by the forward and inverse kernels, including branch and Gamma
denominator guards.

\subsection{Stage 8: special functions}

Each elementary or higher special case has a very small certificate:
\begin{enumerate}
\item a concrete \code{MeijerSpec};
\item its single- or double-residue expansion;
\item a list of Gamma cancellations;
\item a coefficient recurrence matching the target series;
\item a branch equality such as $\sqrt{x^2}=x$;
\item a common tail bound.
\end{enumerate}
For example, the cosine proof compiles to
\begin{lstlisting}[language={}]
SafeCosBranch x p ->
BallEq
  (sqrtPi * meijerG cosSpec (x*x/4) (p+g))
  (cosName x (p+g))
  p
\end{lstlisting}
The induction invariant states equality of the first $N$ rational terms; the
factorial duplication theorem converts the coefficient denominators.  The
Lerch theorem uses the exact identity
\code{poch a k / poch (a+1) k = a / (a+k)} in cross-multiplied form.  Bessel
$Y$ and integer-order $K$ additionally invoke the jet module.

\subsection{Representative object-language target list}

The companion \code{ObjectLanguageTargets.lean} records the following final
families.
\begin{enumerate}
\item \code{HyperTermRecurrence}: exact cleared recurrence for every compiled
      hypergeometric term.
\item \code{HyperTail}: a ratio certificate implies the requested tail ball.
\item \code{SlaterRight} and \code{SlaterLeft}: residue evaluator equals the
      compiled hypergeometric sum.
\item \code{ContourEqualsResidues}: the finite contour evaluator agrees with
      the residue evaluator.
\item \code{MeijerODE}: the two sides of the differential equation agree.
\item \code{PowerShift}, \code{Inversion}, \code{CancelTopBottom}, and
      \code{MultiplyParameters}.
\item Integral targets \code{Mellin}, \code{Euler}, \code{Convolution},
      \code{RationalPower}, and \code{Laplace}, each under an explicit
      certificate predicate.
\item Transform targets \code{KernelPair}, \code{Narain},
      \code{GenLaplace}, and \code{Whittaker} for fixed coded test classes.
\item \code{ExpCase}, \code{CosCase}, \code{LogCase},
      \code{IncompleteGammaCase}, \code{BesselCases}, and \code{LerchCases}.
\end{enumerate}
All equality targets are \code{BallEq} schemas in the requested precision.
Exact finite identities, such as parameter-list normalization, are proved as
ordinary object equalities and then used inside the ball proofs.

\subsection{Milestones and trust audit}

A practical sequence of independently useful milestones is:
\begin{enumerate}
\item checked exact rational/complex balls and hypergeometric recurrence;
\item checked geometric/factorial tail theorem;
\item cosine, exponential, logarithm, and rational special cases entirely in
      \code{Prf PRAAx};
\item generic nonconfluent Slater evaluator;
\item Laurent-jet residues and Bessel $K,Y$;
\item parameter algebra and the differential equation;
\item contour/residue agreement;
\item Mellin, convolution, and Laplace certificates;
\item transform-pair instances.
\end{enumerate}

Each completed module should be audited for \code{sorry}, \code{admit},
\code{axiom}, \code{Classical.choice}, and untrusted native decision procedures.
Use \code{\#print axioms} on host helper theorems, and additionally inspect the
final theorem type to confirm that it is a \code{Prf PRAAx} derivation.  Python
and Wolfram checks in this archive are regression evidence only; they are not
part of the trusted proof base.

\section{Coverage ledger and source audit}\label{sec:coverage}

\subsection{Status vocabulary}

The companion machine-readable ledger contains 107 claim families extracted
from the two requested pages.  The statuses are intentionally stricter than a
binary ``covered/not covered'' flag.

\begin{description}[style=nextline,leftmargin=3.2em]
\item[\statusP\ --- proved]
A complete finite constructive argument is written in this report.  It may use
a previously stated finite Gamma, power-series, or arithmetic lemma, but no
unnamed convergence or choice principle.
\item[\statusQ\ --- quantitatively conditional]
The identity is proved from an explicit certificate predicate: pole clearance,
branch margin, tail majorant, quadrature bound, endpoint modulus, or transform
test-class certificate.  The source page suppresses these hypotheses.
\item[\statusC\ --- corrected]
The literal web statement is too broad, branch-dependent, resonant, or
undefined at a boundary.  A corrected theorem is stated and proved.
\item[\statusR\ --- reduced]
The claim is reduced to the general certified compiler or transform-kernel
principle, but a long specialized normalization or an auxiliary
special-function reduction is not printed in full.  Such a result is not used
as an unproved premise elsewhere.
\item[\statusX\ --- not established]
Reserved for a rejected or unsupported claim.  No ledger item has this status:
problematic source statements were replaced by precise corrected forms rather
than simply discarded.
\end{description}

The distribution is
\[
 46\ \statusP,\qquad 30\ \statusQ,\qquad
 29\ \statusC,\qquad 2\ \statusR,
\]
for a total of 107 source claim families.  Counting a corrected theorem as a
success is justified only because the report both identifies the defect and
proves the repaired statement.

\subsection{Principal corrections}

The most important corrections are mathematical, not stylistic.

\begin{enumerate}
\item The fundamental systems following the differential equation require
nonresonance, or else logarithmic Frobenius solutions and Laurent jets.  Merely
listing $p$ or $q$ power-series solutions does not prove linear independence
when exponents differ by integers.
\item Every occurrence of ``for all $x$'' involving $\Log$, $x^\alpha$,
$\sqrt{x^2}$, inverse trigonometric functions, Bessel functions, or the Lerch
transcendent needs a branch and domain statement.
\item The two step-function formulas are naturally positive-real Mellin
formulas.  For a signed real argument the symmetric statement uses $G(|x|)$,
and the value at $x=1$ depends on the Heaviside convention.
\item The Lerch series formulas are proved first for $|x|<1$; outside the unit
disc the $G$-expression selects a branch continuation, not a globally
single-valued power series.
\item The transform sections use phrases such as ``suitable function'' and
omit convergence restrictions.  The report replaces them by a coded test
class with explicit Mellin, endpoint, Fubini, and inversion certificates.
\item The long six-term representation of $x$ is interpreted at $x=0$ by a
common continuous limit; its individual summands contain singular-looking
factors there.
\item Negative derivative order means an antiderivative only after a base point
or integration constant is specified.  The two displayed G-continuations need
not remain equal for negative order.
\end{enumerate}

\subsection{Levels of verification}

Four levels must not be conflated.

\begin{enumerate}
\item \emph{Mathematical proof.}  The present text supplies finite estimates,
coefficient calculations, and certificate reductions.
\item \emph{Executable regression evidence.}  The Python scripts check exact
coefficient identities and high-precision numerical instances.  Wolfram
Language was used independently to verify the four MathWorld special cases.
\item \emph{Lean host scaffold.}  The files in \code{lean/} define the proposed
records and target interfaces.  They contain placeholders and are explicitly
not presented as a compiled analytic proof.
\item \emph{Object-iPRA theorem.}  The eventual endpoint is a derivation
\code{[] |- [PRAAx] phi}.  No main analytic theorem in this archive is claimed
to have reached that endpoint.
\end{enumerate}

This distinction follows the current ProveIt repository itself: its checked
\code{ConciseIPRA} and \code{IPRADerivations} libraries are kernel-checked,
whereas the proof packets under \code{docs/proofs} are research reports and
formalization plans unless a separate derivation module is listed.

\subsection{Detailed ledger}

The table below is generated from \code{claims/coverage.csv}; the CSV also
contains notes omitted from the typeset table.  In the source column,
\code{WP} abbreviates Wikipedia and \code{MW} abbreviates MathWorld.

\begingroup\scriptsize
\setlength{\tabcolsep}{2pt}
\begin{longtable}{@{}>{\raggedright\arraybackslash}p{0.045\textwidth}>{\raggedright\arraybackslash}p{0.075\textwidth}>{\raggedright\arraybackslash}p{0.12\textwidth}>{\raggedright\arraybackslash}p{0.29\textwidth}>{\centering\arraybackslash}p{0.04\textwidth}>{\raggedright\arraybackslash}p{0.095\textwidth}>{\raggedright\arraybackslash}p{0.245\textwidth}@{}}
\toprule
ID & Source & Topic & Claim & St. & Report & Constructive reading\\
\midrule
\endfirsthead
\toprule
ID & Source & Topic & Claim & St. & Report & Constructive reading\\
\midrule
\endhead
\midrule\multicolumn{7}{r}{Continued on next page}\\
\endfoot
\bottomrule
\endlastfoot
W001 & WP \S{}1 & definition & Mellin--Barnes definition with Gamma(b\_\allowbreak{}j-s), Gamma(1-a\_\allowbreak{}j+s), z\textasciicircum{}\allowbreak{}s. & \statusP & \S{}4.2 & Finite certified contour integral; residue form is primary.\\
W002 & WP\slash{}MW & definition & The s and -s sign conventions are equivalent. & \statusP & \S{}4.2 & Finite substitution u=\allowbreak{}-s on truncated contours.\\
W003 & WP \S{}1 & definition & Index bounds 0$\le$m$\le$q and 0$\le$n$\le$p. & \statusP & \S{}4.1 & Finite list-length predicate.\\
W004 & WP \S{}1 & definition & Numerator pole families must not coincide. & \statusP & \S{}4.1 & Arithmetic progression disjointness plus finite clearance.\\
W005 & WP \S{}1 & definition & The usual definition assumes z$\ne$0. & \statusC & \S{}2.3, \S{}4.5 & Replaced by a quantitative away-from-zero branch certificate.\\
W006 & WP \S{}1 & convergence & Vertical contour converges for |arg z|<$\pi$$\delta$, $\delta$>0. & \statusQ & \S{}5.1 & Explicit exponential tail modulus from Stirling bounds.\\
W007 & WP \S{}1 & convergence & Vertical boundary convergence under the $\sigma$,$\nu$ inequality. & \statusQ & \S{}5.1 & Explicit polynomial tail with rational margin $\eta$>0.\\
W008 & WP \S{}1 & convergence & Right loop converges for q>p. & \statusQ & \S{}5.2 & Factorial\slash{}geometric closing-side bound.\\
W009 & WP \S{}1 & convergence & Right loop converges for p=\allowbreak{}q and |z|<1. & \statusQ & \S{}5.2 & Geometric residue tail with |z|$\le$1-$\epsilon$.\\
W010 & WP \S{}1 & convergence & Right loop boundary |z|=\allowbreak{}1 under Re $\nu$<-1. & \statusQ & \S{}5.2 & p-series tail with explicit margin.\\
W011 & WP \S{}1 & convergence & Left loop converges for p>q. & \statusQ & \S{}5.2 & Inverse-factorial closing-side bound.\\
W012 & WP \S{}1 & convergence & Left loop converges for p=\allowbreak{}q and |z|>1. & \statusQ & \S{}5.2 & Geometric tail in z\textasciicircum{}\allowbreak{}\{-1\}.\\
W013 & WP \S{}1 & convergence & Left loop boundary |z|=\allowbreak{}1 under Re $\nu$<-1. & \statusQ & \S{}5.2 & p-series tail with explicit margin.\\
W014 & WP \S{}1 & contours & Different admissible contours give the same value. & \statusQ & \S{}5.3 & Finite residue rectangles plus vanishing side bounds.\\
W015 & WP \S{}1 & analyticity & G is analytic away from possible singular sets. & \statusQ & \S{}5.5 & Certified local Taylor names and uniform contour bounds.\\
W016 & WP introduction & closure & Fixed G-names are closed under differentiation and antiderivatives. & \statusC & \S{}9 & Fixed-order formulas with branch\slash{}constant data; not an unrestricted class quantifier.\\
W017 & WP \S{}1.1 & ODE & G satisfies the order max(p,q) Meijer differential equation. & \statusP & \S{}6.1 & Kernel polynomial identity or coefficient recurrence.\\
W018 & WP \S{}1.1 & ODE basis & The q displayed functions form a fundamental set when p$\le$q. & \statusC & \S{}6.2 & Proved under nonresonance\slash{}invertible leading coefficient; jets replace resonant powers.\\
W019 & WP \S{}1.1 & ODE basis & The p displayed functions form a fundamental set when p$\ge$q. & \statusC & \S{}6.2 & Proved under dual nonresonance; logarithmic Frobenius blocks in resonant cases.\\
W020 & WP \S{}1.1 & singularities & The displayed basis has only the stated possible singularities. & \statusC & \S{}6.3 & Local certified statement; branch cuts and p=\allowbreak{}q point are explicit.\\
W021 & WP \S{}2 & Slater & Right-loop residue expansion as a sum of pFq$-$1 functions. & \statusP & \S{}5.4 & Explicit residue coefficient and hypergeometric tail.\\
W022 & WP \S{}2 & Slater & Right Slater conditions p<q or p=\allowbreak{}q, |z|<1 and nonconfluence. & \statusP & \S{}5.4 & Exactly the ratio-tail and pole-separation hypotheses.\\
W023 & WP \S{}2 & zero case & For m=\allowbreak{}0 the right-loop value vanishes in its domain. & \statusP & \S{}5.4 & No enclosed poles plus closing-side bound.\\
W024 & WP \S{}2 & Slater & Left-loop residue expansion as a sum of qFp$-$1 functions. & \statusP & \S{}5.4 & Dual residue calculation.\\
W025 & WP \S{}2 & Slater & Left Slater conditions p>q or p=\allowbreak{}q, |z|>1 and nonconfluence. & \statusP & \S{}5.4 & Dual ratio-tail and pole-separation hypotheses.\\
W026 & WP \S{}2 & zero case & For n=\allowbreak{}0 the left-loop value vanishes in its domain. & \statusP & \S{}5.4 & No enclosed poles plus closing-side bound.\\
W027 & WP \S{}2 & hypergeometric & First pFq-to-G representation. & \statusP & \S{}7.1 & One-pole Slater coefficient identity.\\
W028 & WP \S{}2 & hypergeometric & Second pFq-to-G representation at reciprocal argument. & \statusP & \S{}7.1 & Inversion applied to W027.\\
W029 & WP \S{}2 & hypergeometric & Usual pFq convergence regions. & \statusP & \S{}3.3, \S{}7.1 & Primitive-recursive ratio\slash{}factorial tail bounds.\\
W030 & WP \S{}2 & continuation & For p=\allowbreak{}q+1 the G expression continues pFq across |z|=\allowbreak{}1. & \statusC & \S{}7.2 & Branch-specific chain of certified discs\slash{}contours; cut stated explicitly.\\
W031 & WP \S{}2 & extension & G gives a ``natural extension'' for p>q+1. & \statusC & \S{}7.2 & A selected contour-summation name, not a unique branch-free divergent-series sum.\\
W032 & WP \S{}2.1 & polynomial & Terminating hypergeometric polynomial equals the printed sum of two G-functions. & \statusP & \S{}7.3 & Finite residue identity; divergent Gamma prefactors cancel in the pair.\\
W033 & WP \S{}2.1 & polynomial & m,n may be chosen in the stated ranges to avoid singular placements. & \statusP & \S{}7.3 & Finite search over block placements with apartness checks.\\
W034 & WP \S{}2.1 & polynomial & One polynomial-repair G term vanishes in the stated m=\allowbreak{}0 or n=\allowbreak{}0 cases. & \statusP & \S{}7.3 & Empty-pole-family theorem.\\
W035 & WP \S{}3 & cancellation & Cancel equal A$-$ and B+ parameters, lowering p,q,n. & \statusP & \S{}8.1 & Literal cancellation of one Gamma numerator\slash{}denominator factor.\\
W036 & WP \S{}3 & cancellation & Cancel equal A+ and B$-$ parameters, lowering p,q,m. & \statusP & \S{}8.1 & Dual Gamma cancellation.\\
W037 & WP \S{}3 & shift & z\textasciicircum{}\allowbreak{}$\rho$ shifts every a and b parameter by $\rho$. & \statusP & \S{}8.2 & Move z\textasciicircum{}\allowbreak{}$\rho$ into the Mellin kernel on a fixed Log branch.\\
W038 & WP \S{}3 & integer swap & Swap the two separated top parameters when their difference is integral. & \statusP & \S{}8.3 & Finite Gamma recurrence and sign parity.\\
W039 & WP \S{}3 & integer swap & Swap the two separated bottom parameters when their difference is integral. & \statusP & \S{}8.3 & Dual finite Gamma recurrence.\\
W040 & WP \S{}3 & block move & Move an integer-separated top\slash{}bottom pair across numerator blocks. & \statusP & \S{}8.3 & Finite recurrence product and sign.\\
W041 & WP \S{}3 & inversion & G\_\allowbreak{}pq\textasciicircum{}\allowbreak{}\{mn\}(z;a;b)=\allowbreak{}G\_\allowbreak{}qp\textasciicircum{}\allowbreak{}\{nm\}(z\textasciicircum{}\allowbreak{}\{-1\};1-b;1-a). & \statusP & \S{}8.4 & Finite substitution s$\mapsto$-s and list reversal.\\
W042 & WP \S{}3 & multiplication & Integer h multiplication theorem with hp,hq expanded parameters. & \statusP & \S{}8.5 & Gauss Gamma multiplication and explicit constant exponent.\\
W043 & WP \S{}3.1 & derivative & First derivative formula shifting a$_1$ down. & \statusP & \S{}9.1 & Differentiate finite series\slash{}contour and use one Gamma recurrence.\\
W044 & WP \S{}3.1 & derivative & First derivative formula shifting a\_\allowbreak{}p down. & \statusP & \S{}9.1 & Dual block recurrence.\\
W045 & WP \S{}3.1 & derivative & First derivative formula shifting b$_1$ up. & \statusP & \S{}9.1 & Dual lower-block recurrence.\\
W046 & WP \S{}3.1 & derivative & First derivative formula shifting b\_\allowbreak{}q up. & \statusP & \S{}9.1 & Dual lower-block recurrence.\\
W047 & WP \S{}3.1 & derivative & Euler derivative zG$\prime$ expressed with shifted a$_1$ plus (a$_1$$-$1)G. & \statusP & \S{}9.1 & Algebraic expansion of W043.\\
W048 & WP \S{}3.1 & higher derivative & First G representation of z\textasciicircum{}\allowbreak{}h d\textasciicircum{}\allowbreak{}hG\slash{}dz\textasciicircum{}\allowbreak{}h. & \statusP & \S{}9.2 & Falling-factorial polynomial in Mellin variable.\\
W049 & WP \S{}3.1 & higher derivative & Second equivalent G representation of z\textasciicircum{}\allowbreak{}h d\textasciicircum{}\allowbreak{}hG\slash{}dz\textasciicircum{}\allowbreak{}h. & \statusP & \S{}9.2 & Parameter block move applied to W048.\\
W050 & WP \S{}3.1 & reciprocal derivative & Two G representations for derivatives of G(z\textasciicircum{}\allowbreak{}\{-1\}). & \statusP & \S{}9.2 & Chain rule encoded as finite coefficient identity.\\
W051 & WP \S{}3.1 & antiderivative & Negative h gives antiderivative formulas. & \statusC & \S{}9.2 & Each branch includes an integration constant\slash{}base point; the two displayed continuations need not coincide.\\
W052 & WP \S{}3.2 & recurrence & Recurrence for the diagonal pair a$_1$,b$_1$. & \statusP & \S{}9.3 & Subtract two Gamma recurrences in the kernel.\\
W053 & WP \S{}3.2 & recurrence & Recurrence for the diagonal pair a$_1$,b\_\allowbreak{}q. & \statusP & \S{}9.3 & Finite kernel polynomial identity.\\
W054 & WP \S{}3.2 & recurrence & Recurrence for the diagonal pair b$_1$,a\_\allowbreak{}p. & \statusP & \S{}9.3 & Finite kernel polynomial identity.\\
W055 & WP \S{}3.2 & recurrence & Recurrence for the diagonal pair a\_\allowbreak{}p,b\_\allowbreak{}q. & \statusP & \S{}9.3 & Finite kernel polynomial identity.\\
W056 & WP \S{}3.3 & scaling series & Scaling expansion shifting b$_1$ by h. & \statusQ & \S{}9.4 & Taylor theorem at w=\allowbreak{}1 with explicit radius\slash{}tail certificate.\\
W057 & WP \S{}3.3 & scaling series & Scaling expansion shifting b\_\allowbreak{}q by h. & \statusQ & \S{}9.4 & Dual Taylor certificate.\\
W058 & WP \S{}3.3 & scaling series & Scaling expansion shifting a$_1$ by $-$h. & \statusQ & \S{}9.4 & Taylor certificate for z\slash{}w.\\
W059 & WP \S{}3.3 & scaling series & Scaling expansion shifting a\_\allowbreak{}p by $-$h. & \statusQ & \S{}9.4 & Dual Taylor certificate.\\
W060 & WP \S{}4 & Mellin & Mellin transform of G equals its Gamma ratio. & \statusQ & \S{}10.1 & Finite Mellin inversion\slash{}Fej\'er certificate and explicit strip.\\
W061 & WP \S{}4 & Euler integral & Euler integral on (0,1). & \statusQ & \S{}10.2 & Finite Fubini plus Beta contract and endpoint bounds.\\
W062 & WP \S{}4 & Euler integral & Euler integral on (1,$\infty$). & \statusQ & \S{}10.2 & Substitution x=\allowbreak{}1\slash{}t plus Beta contract.\\
W063 & WP \S{}4 & convolution & First concatenated-parameter expression for $\int$G($\eta$x)G($\omega$x)dx. & \statusQ & \S{}10.3 & Mellin transform of second factor at 1$-$s.\\
W064 & WP \S{}4 & convolution & Second equivalent convolution expression. & \statusQ & \S{}10.3 & Reverse the two factors or use inversion.\\
W065 & WP \S{}4.1 & Laplace & Laplace transform of x\textasciicircum{}\allowbreak{}\{-$\alpha$\}G($\eta$x). & \statusQ & \S{}10.5 & Finite Fubini and Gamma integral.\\
W066 & WP \S{}4.1 & Laplace & Bromwich inverse Laplace formula. & \statusQ & \S{}10.5 & Finite Fourier inversion with line\slash{}tail certificate.\\
W067 & WP \S{}4.1 & Laplace & Quadratic-argument Laplace formula. & \statusQ & \S{}10.5 & Gamma duplication after finite Fubini.\\
W068 & WP \S{}5 & transform definition & Definition of a transform-kernel pair and symmetry. & \statusC & \S{}10.6 & Replaced by finite-window inversion on a coded test class.\\
W069 & WP \S{}5.1 & Narain & Narain kernels are an asymmetric transform pair under balance conditions. & \statusQ & \S{}10.6 & Mellin multipliers cancel; endpoint\slash{}test-class certificates added.\\
W070 & WP \S{}5.1 & Narain & Narain pair becomes symmetric under paired-parameter conditions. & \statusQ & \S{}10.6 & Parameter inversion and list normalization.\\
W071 & WP \S{}5.1 & Hankel & Hankel transform is the stated Narain special case. & \statusQ & \S{}10.6, \S{}12.2 & Substitute the Bessel-J G representation.\\
W072 & WP \S{}5.2 & Wimp & Wimp's two kernels form an asymmetric transform pair. & \statusR & \S{}10.6 & Reduced to the certified kernel principle after expanding A; maximal test class not developed.\\
W073 & WP \S{}5.3 & generalized Laplace & Generalized Laplace kernels invert under the balance condition. & \statusQ & \S{}10.6 & Reciprocal Mellin multipliers plus Bromwich certificate.\\
W074 & WP \S{}5.3 & ordinary Laplace & $\gamma$=\allowbreak{}1\slash{}2, $\rho$=\allowbreak{}0, p=\allowbreak{}q=\allowbreak{}0 gives the ordinary Laplace transform. & \statusP & \S{}10.5, \S{}11.1 & Direct exponential special case.\\
W075 & WP \S{}5.4 & Meijer transform & K$\nu$\slash{}I$\nu$ transform pair. & \statusQ & \S{}10.6, \S{}12.2 & Specialize generalized Laplace and substitute Bessel G forms.\\
W076 & WP \S{}5.4 & Meijer transform & Whittaker W\slash{}M transform pair. & \statusQ & \S{}10.6 & Specialization of generalized Laplace plus constructive 1F1\slash{}U-to-G identities.\\
W077 & WP \S{}6 & elementary & e\textasciicircum{}\allowbreak{}x as G\_\allowbreak{}01\textasciicircum{}\allowbreak{}\{10\}(-x|0). & \statusP & \S{}11.1 & Single 0F0 residue series.\\
W078 & WP \S{}6 & elementary & cos x as $\sqrt{\,}$$\pi$ G\_\allowbreak{}02\textasciicircum{}\allowbreak{}\{10\}(x$^2$\slash{}4|0,1\slash{}2). & \statusP & \S{}11.1 & 0F1 coefficient identity and Gamma duplication.\\
W079 & WP \S{}6 & elementary & sin x as $\sqrt{\,}$$\pi$ G\_\allowbreak{}02\textasciicircum{}\allowbreak{}\{10\}(x$^2$\slash{}4|1\slash{}2,0). & \statusC & \S{}11.1 & Proved with explicit $\sqrt{\,}$(x$^2$)=\allowbreak{}x branch sector.\\
W080 & WP \S{}6 & elementary & cosh x as $\sqrt{\,}$$\pi$ G\_\allowbreak{}02\textasciicircum{}\allowbreak{}\{10\}($-$x$^2$\slash{}4|0,1\slash{}2). & \statusP & \S{}11.1 & Coefficient substitution x$\mapsto$ix.\\
W081 & WP \S{}6 & elementary & sinh x as $-$i$\sqrt{\,}$$\pi$ G\_\allowbreak{}02\textasciicircum{}\allowbreak{}\{10\}($-$x$^2$\slash{}4|1\slash{}2,0). & \statusC & \S{}11.1 & Proved with $\sqrt{\,}$($-$x$^2$)=\allowbreak{}ix branch certificate.\\
W082 & WP \S{}6 & elementary & arcsin x as the displayed G\_\allowbreak{}22\textasciicircum{}\allowbreak{}\{12\}. & \statusC & \S{}11.2 & Hypergeometric series plus explicit branch\slash{}cut.\\
W083 & WP \S{}6 & elementary & arctan x as the displayed G\_\allowbreak{}22\textasciicircum{}\allowbreak{}\{12\}. & \statusC & \S{}11.2 & Coefficient series plus explicit branch\slash{}cut.\\
W084 & WP \S{}6 & elementary & arccot x as the displayed G\_\allowbreak{}22\textasciicircum{}\allowbreak{}\{21\}. & \statusC & \S{}11.2 & Inversion; branch normalized as arctan(1\slash{}x).\\
W085 & WP \S{}6 & elementary & ln(1+x) as G\_\allowbreak{}22\textasciicircum{}\allowbreak{}\{12\}. & \statusC & \S{}11.2 & Proved off the branch cut and pole z=\allowbreak{}$-$1.\\
W086 & WP \S{}6 & elementary & Printed six-term G combination equals x. & \statusQ & \S{}11.6 & Confluent jet\slash{}contiguous certificate; full normalized certificate omitted from prose.\\
W087 & WP \S{}6 & step & H(1$-$|x|) as G\_\allowbreak{}11\textasciicircum{}\allowbreak{}\{10\}(x|1;0). & \statusC & \S{}11.5 & Valid on x>0, or use G(|x|); x=\allowbreak{}1 needs a convention.\\
W088 & WP \S{}6 & step & H(|x|$-$1) as G\_\allowbreak{}11\textasciicircum{}\allowbreak{}\{01\}(x|1;0). & \statusC & \S{}11.5 & Same positive-real and boundary correction.\\
W089 & WP \S{}6 & incomplete gamma & Lower incomplete gamma as G\_\allowbreak{}12\textasciicircum{}\allowbreak{}\{11\}. & \statusC & \S{}12.1 & Series\slash{}integral proof on certified $\alpha$,x domain; recurrence continues it.\\
W090 & WP \S{}6 & incomplete gamma & Upper incomplete gamma as G\_\allowbreak{}12\textasciicircum{}\allowbreak{}\{20\}. & \statusC & \S{}12.1 & Two residues; confluent jet at integral parameter differences.\\
W091 & WP \S{}6 & Bessel & J$\nu$ as G\_\allowbreak{}02\textasciicircum{}\allowbreak{}\{10\}. & \statusC & \S{}12.2 & 0F1 series with explicit square-root branch.\\
W092 & WP \S{}6 & Bessel & Y$\nu$ as G\_\allowbreak{}13\textasciicircum{}\allowbreak{}\{20\}. & \statusC & \S{}12.2 & Two residues reduce to cot$\pi$$\nu$ J$\nu$$-$csc$\pi$$\nu$ J$-$$\nu$; jets at integer $\nu$.\\
W093 & WP \S{}6 & Bessel & I$\nu$ as i\textasciicircum{}\allowbreak{}\{-$\nu$\}G\_\allowbreak{}02\textasciicircum{}\allowbreak{}\{10\}($-$x$^2$\slash{}4). & \statusC & \S{}12.2 & J$\nu$(ix) with branch certificate.\\
W094 & WP \S{}6 & Bessel & K$\nu$ as 1\slash{}2 G\_\allowbreak{}02\textasciicircum{}\allowbreak{}\{20\}. & \statusC & \S{}12.2 & Two 0F1 residues; Laurent jets at integer $\nu$.\\
W095 & WP \S{}6 & Lerch & $\Phi$(x,n,a) as repeated-parameter G for n$\ge$0. & \statusC & \S{}12.3 & Coefficient identity for |x|<1; branch continuation otherwise.\\
W096 & WP \S{}6 & Lerch & $\Phi$(x,$-$n,a) as repeated-parameter G for n$\ge$0. & \statusC & \S{}12.3 & Coefficient identity for |x|<1; branch continuation otherwise.\\
W097 & WP \S{}6 & parameter derivative & Fixed $\alpha$-derivatives of lower\slash{}upper incomplete gamma are expressible with G. & \statusP & \S{}12.4 & Explicit finite hypergeometric\slash{}G sum for every fixed derivative order.\\
M001 & MW (1)--(3) & definition & Minus-sign and plus-sign Mellin--Barnes definitions. & \statusP & \S{}4.2 & Finite substitution proves equivalence.\\
M002 & MW (3) & generalized G & Generalized fourth parameter r. & \statusP & \S{}4.4 & Certified rational r and branch; substitution s=\allowbreak{}ru.\\
M003 & MW (4) & special case & G\_\allowbreak{}22\textasciicircum{}\allowbreak{}\{12\}(z|1,1;1,0)=\allowbreak{}log(1+z). & \statusC & \S{}11.2 & Domain\slash{}cut restored; direct 2F1 series.\\
M004 & MW (5) & special case & G\_\allowbreak{}22\textasciicircum{}\allowbreak{}\{12\}(z|1,1;1,1)=\allowbreak{}z\slash{}(1+z). & \statusP & \S{}11.2 & Hypergeometric parameter cancellation.\\
M005 & MW (6) & special case & G\_\allowbreak{}02\textasciicircum{}\allowbreak{}\{10\}(z\slash{}2|0,1\slash{}2)=\allowbreak{}cos($\sqrt{\,}$(2z))\slash{}$\sqrt{\,}$$\pi$. & \statusC & \S{}11.3 & Cosine formula with a selected square-root branch.\\
M006 & MW (7) & special case & G\_\allowbreak{}10\textasciicircum{}\allowbreak{}\{01\}(z|1$-$a)=\allowbreak{}z\textasciicircum{}\allowbreak{}\{-a\}e\textasciicircum{}\allowbreak{}\{-1\slash{}z\}. & \statusC & \S{}11.3 & Single left residue with Log branch.\\
M007 & MW (8) & generalized special case & Generalized-r formula at r=\allowbreak{}1\slash{}2 equals cos z\slash{}$\sqrt{\,}$$\pi$. & \statusP & \S{}11.3 & Direct residue coefficient identity.\\
B001 & WP introduction & scope & Most established special functions can be represented as G-functions. & \statusC & \S{}7, \S{}12.5 & Interpreted as closure under a specified finite grammar of hypergeometric\slash{}G operations.\\
B002 & WP introduction & products\slash{}integrals & Products G$_1$(cx\textasciicircum{}\allowbreak{}$\gamma$)G$_2$(dx\textasciicircum{}\allowbreak{}$\delta$) with rational $\gamma$\slash{}$\delta$ integrate to another G. & \statusQ & \S{}10.4 & Rational-power Mellin compiler using s=\allowbreak{}B u and Gauss multiplication.\\
B003 & WP introduction & numerics & Numerical contour integration is a practical evaluation method. & \statusR & \S{}4--\S{}5, companion code & A certified quadrature design is supplied; performance is not a mathematical theorem.\\
\end{longtable}
\endgroup


\appendix
\section{Explicit estimates and certificate templates}\label{sec:appendices}

This appendix records details that are routine in classical notation but
important for an iPRA implementation: every ``eventually'', ``uniformly'', and
``interchange the order'' is replaced by a bounded arithmetic inequality.

\subsection{A uniform hypergeometric ratio bound}

Let
\[
 T_k=\frac{\prod_{i=1}^{r}(\alpha_i)_k}
           {\prod_{j=1}^{s}(\beta_j)_k}\frac{z^k}{k!}.
\]
Suppose rational bounds $\abs{\alpha_i}\le B$ and
$\abs{\beta_j}\le B$ are certified and every $\beta_j+k$ is apart from zero.
For $k\ge2B+1$,
\[
 \abs{k+\alpha_i}\le\frac32k,
 \qquad
 \abs{k+\beta_j}\ge\frac12k,
 \qquad k+1\ge k.
\]
Consequently
\begin{equation}\label{eq:appendix-ratio}
 \abs{\frac{T_{k+1}}{T_k}}
 \le \abs z\,\frac{3^r2^{s-r}}{k^{s+1-r}}.
\end{equation}
For $r\le s$, bounded search finds $N$ making the right side at most $1/2$;
then
\begin{equation}\label{eq:appendix-tail-half}
 \abs{\sum_{k=N}^{\infty}T_k}\le2\abs{T_N}.
\end{equation}
For $r=s+1$ and $\abs z\le1-\epsilon$, use the sharper expansion
\[
 \frac{\prod_i(k+\alpha_i)}{(k+1)\prod_j(k+\beta_j)}
 =1+O(1/k)
\]
with a rational constant $C$.  Choosing
$N\ge2C/\epsilon$ gives ratio at most $1-\epsilon/2$ and tail
$2\abs{T_N}/\epsilon$.  The $O(1/k)$ assertion is implemented by expanding the
two monic integer-cleared polynomials and bounding their lower coefficients.
No asymptotic notation appears in the checker.

For $r>s+1$, the nonterminating power series generally has zero radius.  The
report never calls this series convergent; the corresponding $G$-function is a
contour/residue summation name.  If an upper parameter is a nonpositive integer,
the series terminates and is a finite polynomial, covered by
\cref{thm:polynomial-repair}.

\subsection{Derivation of the right Slater coefficient}

Put $s=b_h+k+u$.  Since
\[
 \Gamma(-k-u)=\frac{(-1)^{k+1}}{k!u}+O(1),
\]
the ordinary residue of the integrand has a leading minus sign.  The right
loop is negatively oriented, so the contour integral contributes the positive
coefficient $(-1)^k/k!$.  The remaining factors are rewritten using
\begin{align*}
 \Gamma(c+k)&=\Gamma(c)(c)_k,\\
 \Gamma(c-k)&=(-1)^k\frac{\Gamma(c)}{(1-c)_k},\\
 \frac1{\Gamma(c-k)}&=(-1)^k\frac{(1-c)_k}{\Gamma(c)}.
\end{align*}
After collecting signs, the argument is
$(-1)^{p-m-n}z$ and the constant term is
\begin{equation}\label{eq:appendix-Ch}
 C_h=
 \frac{\displaystyle
   \prod_{\substack{1\le j\le m\\j\ne h}}\Gamma(b_j-b_h)
   \prod_{j=1}^{n}\Gamma(1+b_h-a_j)}
 {\displaystyle
   \prod_{j=m+1}^{q}\Gamma(1+b_h-b_j)
   \prod_{j=n+1}^{p}\Gamma(a_j-b_h)}.
\end{equation}
The $k$th coefficient is therefore
\[
 C_hz^{b_h}
 \frac{\prod_{j=1}^{p}(1+b_h-a_j)_k}
      {\prod_{\substack{1\le j\le q\\j\ne h}}
       (1+b_h-b_j)_k}
 \frac{[(-1)^{p-m-n}z]^k}{k!},
\]
which is exactly the $h$th hypergeometric term in
\cref{thm:slater-right}.  Every equation above has a denominator-cleared
factorial proof for rational parameters away from poles; for general certified
complex parameters it is a finite Gamma recurrence ball equality.

The left formula follows by setting $s=a_h-1-k+u$.  Here the left loop has
positive orientation, and the same three Gamma identities produce the dual
${}_qF_{p-1}$ expression.

\subsection{Confluent residues without parameter limits}

Suppose $r$ Gamma factors have a pole at $s_0$.  Write
\[
 F(s_0+u)=\sum_{j=-r}^{r-1}c_ju^j+u^rR(u),
 \qquad \abs{R(u)}\le M
 \quad(\abs u\le\rho).
\]
Multiplying by
\[
 z^{s_0+u}=z^{s_0}\sum_{j=0}^{r-1}
 \frac{(u\Log z)^j}{j!}+u^rR_z(u)
\]
shows that the residue is
\begin{equation}\label{eq:appendix-confluent-res}
 z^{s_0}\sum_{j=0}^{r-1}c_{-1-j}\frac{(\Log z)^j}{j!}.
\end{equation}
Thus the degree in $\Log z$ is at most $r-1$.  To certify the remainder, use
truncated convolution.  If
$A=A_{<r}+u^rR_A$ and $B=B_{<r}+u^rR_B$ on $\abs u\le\rho$, then
\[
 \abs{R_{AB}}
 \le \abs{A_{<r}}_\rho M_B+
      \abs{B_{<r}}_\rho M_A+\rho^rM_AM_B,
\]
where $\abs{\sum a_ju^j}_\rho=\sum\abs{a_j}\rho^j$ after shifting away
negative powers.  This formula is a primitive-recursive update of rational
bounds.  It is the formal replacement for phrases such as ``take the limit as
$b_2\to b_1$''.

\subsection{The differential equation as a polynomial identity}

Let $\thetaop=z\frac{\dd}{\dd z}$.  Acting on $z^s$ replaces $\thetaop$ by
$s$.  Therefore the differential operator in \cref{thm:ode} multiplies
the kernel by
\[
 P(s)=(-1)^{p-m-n}z\prod_{j=1}^{p}(s-a_j+1)
      -\prod_{j=1}^{q}(s-b_j).
\]
Gamma recurrence gives the exact shift identity
\begin{equation}\label{eq:appendix-kernel-shift}
 \Phi(s-1)=(-1)^{p-m-n}
 \frac{\prod_{j=1}^{p}(s-a_j)}
      {\prod_{j=1}^{q}(s-b_j-1)}\Phi(s),
\end{equation}
with the convention of \eqref{eq:main-definition}.  After replacing $s$ by
$s+1$ and multiplying by $z^{s+1}$, the first term in $P(s)$ is the shifted
second term.  On a finite contour the substitution is exact; the two short
connector segments cancel or tend below the allocated error by the same
vertical estimate.  In a residue series, \eqref{eq:appendix-kernel-shift} is
precisely the coefficient recurrence.  Hence the ODE is proved without a
classical theorem about differentiating arbitrary improper integrals.

\subsection{Why the nonresonance correction is necessary}

For $p\le q$, the $h$th proposed local solution begins with
$C_hz^{b_h}(1+O(z))$.  The first $q$ Euler derivatives at a positive rational
base point have leading matrix
\[
 V=(b_h^{\ell})_{0\le\ell<q,\ 1\le h\le q}.
\]
Its determinant is the Vandermonde product
$\prod_{i<j}(b_j-b_i)$.  A quantitative lower bound for this product proves
linear independence and supplies a conditioning constant.  If two exponents
coincide, or differ by an integer in a regular-singular equation, this argument
fails and logarithmic solutions may occur.  The Laurent-jet construction
produces those solutions.  Therefore the unqualified word ``fundamental'' on
the source page cannot be accepted literally in resonant cases.

\subsection{A finite Fubini budget}

Consider a formally desired exchange
\[
 \int_0^\infty\int_{-\infty}^{\infty}F(x,t)\dd t\dd x
 =\int_{-\infty}^{\infty}\int_0^\infty F(x,t)\dd x\dd t.
\]
At precision $2^{-p}$ choose rational $0<\epsilon<R$ and $T>0$ so that
\begin{align*}
 \int_0^\epsilon\int_{\R}\abs F&\le2^{-p-4},&
 \int_R^\infty\int_{\R}\abs F&\le2^{-p-4},\\
 \int_\epsilon^R\int_{\abs t>T}\abs F&\le2^{-p-4},&
 E_{\rm quad}&\le2^{-p-4}.
\end{align*}
On the finite rectangle $[\epsilon,R]\times[-T,T]$, both iterated integrals
are represented by the same finite double quadrature sum, so reordering is
literal associativity/commutativity.  The four discarded errors total at most
$2^{-p-2}$; the remaining budget covers the inner Gamma/Beta evaluation and
rounding.  This is the pattern used in every integral theorem of
\cref{sec:transforms}.

\subsection{Branch transport by overlapping boxes}

A branch certificate is not merely a symbolic choice of $\Arg$.  It contains a
rational box $D$, a rational center $z_0$, a value $L_0$ approximating
$\Log z_0$, and a margin $d>0$ separating $D$ from zero and the cut.  If
$z_1$ lies in the same box and
$\abs{(z_1-z_0)/(z_1+z_0)}\le r<1$, then
\begin{equation}\label{eq:appendix-log-step}
 \Log z_1-\Log z_0
 =2\sum_{k=0}^{N-1}\frac{w^{2k+1}}{2k+1}+R_N,
 \quad w=\frac{z_1-z_0}{z_1+z_0},
\end{equation}
with
\[
 \abs{R_N}\le\frac{2r^{2N+1}}{(2N+1)(1-r^2)}.
\]
A finite chain of overlapping boxes transports the branch.  Powers are then
$z^a=\exp(a\Log z)$ with a propagated ball.  This construction proves every
branch continuation used in the report and makes monodromy visible: a loop
whose certified logarithm increment is $2\pi\ii$ changes $z^a$ by
$\e^{2\pi\ii a}$ rather than pretending the function is globally
single-valued.

\subsection{Representative iPRA sentence schemas}

For a fixed compiled cosine specification, the mathematical statement becomes
\[
 \forall x\,\forall p\,
 \bigl(\Safe_{\cos}(x,p)\to
 \BallEq(\sqrt\pi\,G(x^2/4),\cos(x),p)\bigr).
\]
The displayed quantifiers are ordinary first-order quantifiers over natural
codes.  The functions \code{SafeCos}, \code{meijerCos}, \code{cosEval}, and
\code{ballEq} are fixed primitive-recursive symbols.  The actual Lean target is
a closed formula built from their graph equations and a proof
\code{[] |- [PRAAx] phi}.

For Mellin transformation, a fixed certificate producer $C$ gives
\[
 \forall u\,\forall p\,
 \bigl(\operatorname{MellinCert}(C,u,p)\to
 \BallEq(\operatorname{MellinEval}(u,p),
         \operatorname{GammaRatio}(u,p),p)\bigr).
\]
This is not second-order quantification over functions or certificates.  The
programs are compiled into the language before the theorem is stated.

\subsection{Reproducibility and limitations}

The archive contains:
\begin{enumerate}
\item the full LaTeX source and compiled PDF;
\item a CSV ledger and its generated typeset table;
\item a Lean-oriented host scaffold and explicit object-language targets;
\item exact coefficient checks using rational arithmetic;
\item high-precision numerical regression checks using \code{mpmath};
\item source provenance, build logs, rendered PDF page images, and checksums.
\end{enumerate}

The executable checks do not prove the universal analytic theorems.  Conversely,
the human proofs do not imply that a Lean file has been compiled.  The missing
formal work is precisely the construction of the object-language derivations
described in \cref{sec:formalization}.  This boundary is part of the result,
not a hidden caveat.


\begin{thebibliography}{99}
\bibitem{DLMF16}
NIST Digital Library of Mathematical Functions, Chapter 16, especially
Sections 16.17--16.21, \emph{Meijer $G$-Function},
\url{https://dlmf.nist.gov/16}, inspected 16 November 2026.

\bibitem{Wikipedia}
Wikipedia contributors, \emph{Meijer $G$-function},
\url{https://en.wikipedia.org/wiki/Meijer_G-function}, inspected 16 November 2026.

\bibitem{MathWorld}
E. W. Weisstein, \emph{Meijer $G$-Function}, MathWorld---A Wolfram Web Resource,
\url{https://mathworld.wolfram.com/MeijerG-Function.html}, inspected 16 November 2026.

\bibitem{Bateman}
A. Erd\'elyi, W. Magnus, F. Oberhettinger, and F. G. Tricomi,
\emph{Higher Transcendental Functions}, Vol. I, McGraw--Hill, 1953,
Section 5.3.

\bibitem{Luke}
Y. L. Luke, \emph{The Special Functions and Their Approximations}, Vol. I,
Academic Press, 1969.

\bibitem{Prudnikov}
A. P. Prudnikov, Yu. A. Brychkov, and O. I. Marichev,
\emph{Integrals and Series}, Vol. 3: More Special Functions, Gordon and Breach,
1990.

\bibitem{Beals}
R. Beals and J. Szmigielski, \emph{Meijer $G$-functions: a gentle
introduction}, Notices of the AMS 60 (2013), 866--872.

\bibitem{ProveIt}
V. Reshetnikov, \emph{ProveIt}, commit
\code{85afffc93b2abbc1972fdfc98244052cbfbda1b3},
\url{https://github.com/VladimirReshetnikov/ProveIt}.

\bibitem{ConciseIPRA}
ProveIt, \path{Logic/PrimitiveRecursiveArithmetic/ConciseIPRA}, checked Lean
formalization of iPRA and the native scalar core, commit as in \cite{ProveIt}.

\bibitem{GammaPacket}
ProveIt, \emph{Gamma in intuitionistic primitive recursive arithmetic},
\path{Logic/PrimitiveRecursiveArithmetic/docs/proofs/analysis/gamma_ipra},
commit as in \cite{ProveIt}.

\bibitem{ResiduePacket}
ProveIt, \emph{Cauchy and residue theorems in iPRA},
\path{Logic/PrimitiveRecursiveArithmetic/docs/proofs/analysis/cauchy_residue_ipra},
commit as in \cite{ProveIt}.

\bibitem{BarnesPacket}
ProveIt, \emph{The Barnes $G$-function in iPRA},
\path{Logic/PrimitiveRecursiveArithmetic/docs/proofs/analysis/barnes_g_ipra_2026-11-16},
commit as in \cite{ProveIt}.
\end{thebibliography}

\end{document}
